
\Data\ID{1003023104}
\Updated{2021-06-21 09:06:00}
\Level{A}
\SubArea{Angles, arcs and sectors}
\LANGen{
\Question{
Select such a pair of numbers that represents measures of two coterminal angles. (Two angles are coterminal if when they are drawn in the standard position their terminal sides merge.)}
{\( \frac{15}2\pi;\ -8.5\pi \)}{\( -21.5\pi;\ -\frac92\pi \)}{\( \frac{19}2\pi;\ 20.5\pi \)}{\( \frac{35}2\pi;\ -\frac{35}2\pi \)}{}{}}
\LANGpl{
\Question{
Wybierz parę liczb, które przedstawiają tą samą miarę kątów.}
{\( \frac{15}2\pi;\ -8{,}5\pi \)}{\( -21{,}5\pi;\ -\frac92\pi \)}{\( \frac{19}2\pi;\ 20{,}5\pi \)}{\( \frac{35}2\pi;\ -\frac{35}2\pi \)}{}{}}
\LANGcs{
\Question{
Vyber takovou dvojici čísel, které představují velikost stejného orientovaného úhlu.}
{\( \frac{15}2\pi;\ -8{,}5\pi \)}{\( -21{,}5\pi;\ -\frac92\pi \)}{\( \frac{19}2\pi;\ 20{,}5\pi \)}{\( \frac{35}2\pi;\ -\frac{35}2\pi \)}{}{}}
\LANGsk{
\Question{
Vyberte tú dvojicu čísel, ktorá predstavuje veľkosť jedného a toho istého orientovaného uhla.}
{\( \frac{15}2\pi;\ -8{,}5\pi \)}{\( -21{,}5\pi;\ -\frac92\pi \)}{\( \frac{19}2\pi;\ 20{,}5\pi \)}{\( \frac{35}2\pi;\ -\frac{35}2\pi \)}{}{}}
\LANGes{
\Question{
Elige la pareja de números que representa medidas de dos ángulos coterminales. (Dos ángulos coterminales son tales que si los dibujamos en la posición estándar, sus lados terminales coinciden.)}
{\( \frac{15}2\pi;\ -8.5\pi \)}{\( -21.5\pi;\ -\frac92\pi \)}{\( \frac{19}2\pi;\ 20.5\pi \)}{\( \frac{35}2\pi;\ -\frac{35}2\pi \)}{}{}}
\End

\Data\ID{1003023105}
\Updated{2021-06-21 09:06:24}
\Level{A}
\SubArea{Angles, arcs and sectors}
\LANGen{
\Question{
Select such a pair of numbers that represents measures of two coterminal angles. (Two angles are coterminal if when they are drawn in the standard position their terminal sides merge.)}
{\( -21.5\pi;\ -41.5\pi \)}{\( -21.5\pi;\ 21.5\pi \)}{\( -21.5\pi; 41.5\pi \)}{\( 21.5\pi;\ -41.5\pi \)}{}{}}
\LANGpl{
\Question{
Wybierz parę liczb, które przedstawiają tą samą miarę kątów.}
{\( -21{,}5\pi;\ -41{,}5\pi \)}{\( -21{,}5\pi;\ 21{,}5\pi \)}{\( -21{,}5\pi; 41{,}5\pi \)}{\( 21{,}5\pi;\ -41{,}5\pi \)}{}{}}
\LANGcs{
\Question{
Vyber takovou dvojici čísel, které představují velikost stejného orientovaného úhlu.}
{\( -21{,}5\pi;\ -41{,}5\pi \)}{\( -21{,}5\pi;\ 21{,}5\pi \)}{\( -21{,}5\pi; 41{,}5\pi \)}{\( 21{,}5\pi;\ -41{,}5\pi \)}{}{}}
\LANGsk{
\Question{
Vyberte tú dvojicu čísel, ktorá predstavuje veľkosť jedného a toho istého orientovaného uhla.}
{\( -21{,}5\pi;\ -41{,}5\pi \)}{\( -21{,}5\pi;\ 21{,}5\pi \)}{\( -21{,}5\pi; 41{,}5\pi \)}{\( 21{,}5\pi;\ -41{,}5\pi \)}{}{}}
\LANGes{
\Question{
Elige la pareja de números que representa medidas de dos ángulos coterminales. (Dos ángulos coterminales son tales que si los dibujamos en la posición estándar, sus lados terminales coinciden.)}
{\( -21.5\pi;\ -41.5\pi \)}{\( -21.5\pi;\ 21.5\pi \)}{\( -21.5\pi; 41.5\pi \)}{\( 21.5\pi;\ -41.5\pi \)}{}{}}
\End

\Data\ID{1003023108}
\Updated{2021-06-21 09:06:39}
\Level{A}
\SubArea{Angles, arcs and sectors}
\LANGen{
\Question{
Select a pair of angles \( \alpha \), \( \beta \), that take the same position on the unit circle?}
{\( \alpha = \frac23\pi;\ \beta=-\frac43\pi \)}{\( \alpha = \frac65\pi;\ \beta=\frac45\pi \)}{\( \alpha = -\frac76\pi;\ \beta = -\frac{32}6\pi \)}{\( \alpha=-\frac23\pi;\ \beta=\frac{31}3\pi \)}{}{}}
\LANGpl{
\Question{
Wybierz parę kątów \( \alpha \), \( \beta \), które  przyjmują taką samą pozycję na okręgu jednostkowym.}
{\( \alpha = \frac23\pi;\ \beta=-\frac43\pi \)}{\( \alpha = \frac65\pi;\ \beta=\frac45\pi \)}{\( \alpha = -\frac76\pi;\ \beta = -\frac{32}6\pi \)}{\( \alpha=-\frac23\pi;\ \beta=\frac{31}3\pi \)}{}{}}
\LANGcs{
\Question{
Která dvojice úhlů \( \alpha \), \( \beta \) má stejné grafické znázornění na jednotkové kružnici?}
{\( \alpha = \frac23\pi;\ \beta=-\frac43\pi \)}{\( \alpha = \frac65\pi;\ \beta=\frac45\pi \)}{\( \alpha = -\frac76\pi;\ \beta = -\frac{32}6\pi \)}{\( \alpha=-\frac23\pi;\ \beta=\frac{31}3\pi \)}{}{}}
\LANGsk{
\Question{
Ktorá dvojica uhlov \( \alpha \), \( \beta \) má rovnaké grafické znázornenie na jednotkovej kružnici?}
{\( \alpha = \frac23\pi;\ \beta=-\frac43\pi \)}{\( \alpha = \frac65\pi;\ \beta=\frac45\pi \)}{\( \alpha = -\frac76\pi;\ \beta = -\frac{32}6\pi \)}{\( \alpha=-\frac23\pi;\ \beta=\frac{31}3\pi \)}{}{}}
\LANGes{
\Question{
Elije la pareja de ángulos \( \alpha \), \( \beta \), que tienen la misma posición en la circunferencia goniométrica.}
{\( \alpha = \frac23\pi;\ \beta=-\frac43\pi \)}{\( \alpha = \frac65\pi;\ \beta=\frac45\pi \)}{\( \alpha = -\frac76\pi;\ \beta = -\frac{32}6\pi \)}{\( \alpha=-\frac23\pi;\ \beta=\frac{31}3\pi \)}{}{}}
\End

\Data\ID{1003023109}
\Updated{2021-06-21 09:06:56}
\Level{A}
\SubArea{Angles, arcs and sectors}
\LANGen{
\Question{
Select such a pair of angles \( \alpha \), \( \beta \) that do not take the same position on the unit circle.}
{\( \alpha= 129^{\circ};\ \beta=859^{\circ} \)}{\( \alpha= 575^{\circ};\ \beta=2015^{\circ} \)}{\( \alpha= \frac23\pi;\ \beta=\frac{14}3\pi \)}{\( \alpha= \frac54\pi;\ \beta=\frac{29}4\pi \)}{}{}}
\LANGpl{
\Question{
Wybierz parę kątów \( \alpha \), \( \beta \), które nie przyjmują takiej samej pozycji na okręgu jednostkowym.}
{\( \alpha= 129^{\circ};\ \beta=859^{\circ} \)}{\( \alpha= 575^{\circ};\ \beta=2015^{\circ} \)}{\( \alpha= \frac23\pi;\ \beta=\frac{14}3\pi \)}{\( \alpha= \frac54\pi;\ \beta=\frac{29}4\pi \)}{}{}}
\LANGcs{
\Question{
Která dvojice úhlů \( \alpha \), \( \beta \) se nezobrazí na jednotkové kružnici stejně?}
{\( \alpha= 129^{\circ};\ \beta=859^{\circ} \)}{\( \alpha= 575^{\circ};\ \beta=2015^{\circ} \)}{\( \alpha= \frac23\pi;\ \beta=\frac{14}3\pi \)}{\( \alpha= \frac54\pi;\ \beta=\frac{29}4\pi \)}{}{}}
\LANGsk{
\Question{
Vyberte tú dvojicu uhlov  \( \alpha \), \( \beta \),  ktoré sa na jednotkovej kružnici nezobrazia rovnako.}
{\( \alpha= 129^{\circ};\ \beta=859^{\circ} \)}{\( \alpha= 575^{\circ};\ \beta=2015^{\circ} \)}{\( \alpha= \frac23\pi;\ \beta=\frac{14}3\pi \)}{\( \alpha= \frac54\pi;\ \beta=\frac{29}4\pi \)}{}{}}
\LANGes{
\Question{
Elige la pareja de ángulos \( \alpha \), \( \beta \), que no tienen la misma posición en la circunferencia goniométrica.}
{\( \alpha= 129^{\circ};\ \beta=859^{\circ} \)}{\( \alpha= 575^{\circ};\ \beta=2015^{\circ} \)}{\( \alpha= \frac23\pi;\ \beta=\frac{14}3\pi \)}{\( \alpha= \frac54\pi;\ \beta=\frac{29}4\pi \)}{}{}}
\End

\Data\ID{1003023110}
\Updated{2021-06-21 09:06:54}
\Level{A}
\SubArea{Angles, arcs and sectors}
\LANGen{
\Question{
Which of the given sets contains measures of three angles that take the same position as the angle \( \varphi = 18^{\circ} \) on the unit circle?}
{\( \left\{ 378^{\circ};\ -342^{\circ};\ -1422^{\circ} \right\} \)}{\( \left\{ 1098^{\circ};\ 1818^{\circ};\ -1052^{\circ} \right\} \)}{\( \left\{ 1098^{\circ};\ -1062^{\circ};\ -1812^{\circ} \right\} \)}{\( \left\{ 378^{\circ};\ 1092^{\circ};\ -1062^{\circ} \right\} \)}{}{}}
\LANGpl{
\Question{
Który z podanych zbiorów zawiera miary trzech kątów przyjmujących taką sama pozycję co kąt   \( \varphi = 18^{\circ} \) na okręgu jednostkowym?}
{\( \left\{ 378^{\circ};\ -342^{\circ};\ -1422^{\circ} \right\} \)}{\( \left\{ 1098^{\circ};\ 1818^{\circ};\ -1052^{\circ} \right\} \)}{\( \left\{ 1098^{\circ};\ -1062^{\circ};\ -1812^{\circ} \right\} \)}{\( \left\{ 378^{\circ};\ 1092^{\circ};\ -1062^{\circ} \right\} \)}{}{}}
\LANGcs{
\Question{
Která množina obsahuje velikosti tří úhlů, které se na jednotkové kružnici zobrazí stejně jako úhel \( \varphi = 18^{\circ} \)?}
{\( \left\{ 378^{\circ};\ -342^{\circ};\ -1422^{\circ} \right\} \)}{\( \left\{ 1098^{\circ};\ 1818^{\circ};\ -1052^{\circ} \right\} \)}{\( \left\{ 1098^{\circ};\ -1062^{\circ};\ -1812^{\circ} \right\} \)}{\( \left\{ 378^{\circ};\ 1092^{\circ};\ -1062^{\circ} \right\} \)}{}{}}
\LANGsk{
\Question{
Ktorá množina obsahuje veľkosti troch uhlov, ktoré sa na jednotkovej kružnici zobrazia rovnako ako uhol   \( \varphi = 18^{\circ} \)?}
{\( \left\{ 378^{\circ};\ -342^{\circ};\ -1422^{\circ} \right\} \)}{\( \left\{ 1098^{\circ};\ 1818^{\circ};\ -1052^{\circ} \right\} \)}{\( \left\{ 1098^{\circ};\ -1062^{\circ};\ -1812^{\circ} \right\} \)}{\( \left\{ 378^{\circ};\ 1092^{\circ};\ -1062^{\circ} \right\} \)}{}{}}
\LANGes{
\Question{
¿Cuál de los conjuntos contiene las medidas de tres ángulos que tienen la misma posición en la circunferencia goniométrica que el ángulo \( \varphi = 18^{\circ} \)?}
{\( \left\{ 378^{\circ};\ -342^{\circ};\ -1422^{\circ} \right\} \)}{\( \left\{ 1098^{\circ};\ 1818^{\circ};\ -1052^{\circ} \right\} \)}{\( \left\{ 1098^{\circ};\ -1062^{\circ};\ -1812^{\circ} \right\} \)}{\( \left\{ 378^{\circ};\ 1092^{\circ};\ -1062^{\circ} \right\} \)}{}{}}
\End

\Data\ID{1003023111}
\Updated{2021-06-21 09:06:43}
\Level{A}
\SubArea{Angles, arcs and sectors}
\LANGen{
\Question{
Select the set that contains measures of two angles that do not take the same position as the angle \( \alpha=\frac3{10}\pi \) radians on the unit circle.}
{\( \left\{ \frac{37}{10}\pi;\ -\frac{37}{10}\pi \right\} \)}{\( \left\{ \frac{63}{10}\pi;\ -\frac{103}{10}\pi \right\} \)}{\( \left\{ -\frac{17}{10}\pi;\ -\frac{77}{10}\pi \right\} \)}{\( \left\{ -\frac{37}{10}\pi;\ \frac{23}{10}\pi \right\} \)}{}{}}
\LANGpl{
\Question{
Który z podanych zbiorów zawiera miary dwóch kątów, które nie przyjmują takiej samej pozycji co kąt \( \alpha=\frac3{10}\pi \) radianów na okręgu jednostkowym.}
{\( \left\{ \frac{37}{10}\pi;\ -\frac{37}{10}\pi \right\} \)}{\( \left\{ \frac{63}{10}\pi;\ -\frac{103}{10}\pi \right\} \)}{\( \left\{ -\frac{17}{10}\pi;\ -\frac{77}{10}\pi \right\} \)}{\( \left\{ -\frac{37}{10}\pi;\ \frac{23}{10}\pi \right\} \)}{}{}}
\LANGcs{
\Question{
Vyberte množinu, která obsahuje velikosti dvou úhlů, jež se na jednotkové kružnici nezobrazí stejně jako úhel \( \alpha=\frac3{10}\pi \) .}
{\( \left\{ \frac{37}{10}\pi;\ -\frac{37}{10}\pi \right\} \)}{\( \left\{ \frac{63}{10}\pi;\ -\frac{103}{10}\pi \right\} \)}{\( \left\{ -\frac{17}{10}\pi;\ -\frac{77}{10}\pi \right\} \)}{\( \left\{ -\frac{37}{10}\pi;\ \frac{23}{10}\pi \right\} \)}{}{}}
\LANGsk{
\Question{
Vyberte množinu, ktorá obsahuje veľkosti dvoch uhlov, ktoré sa nezobrazia rovnako na jednotkovej kružnici ako uhol  \( \alpha=\frac3{10}\pi \) .}
{\( \left\{ \frac{37}{10}\pi;\ -\frac{37}{10}\pi \right\} \)}{\( \left\{ \frac{63}{10}\pi;\ -\frac{103}{10}\pi \right\} \)}{\( \left\{ -\frac{17}{10}\pi;\ -\frac{77}{10}\pi \right\} \)}{\( \left\{ -\frac{37}{10}\pi;\ \frac{23}{10}\pi \right\} \)}{}{}}
\LANGes{
\Question{
Elige el conjunto que contiene las medidas de dos ángulos que no tienen la misma posición en la circunferencia goniométrica que el ángulo \( \alpha=\frac3{10}\pi \) radianes.}
{\( \left\{ \frac{37}{10}\pi;\ -\frac{37}{10}\pi \right\} \)}{\( \left\{ \frac{63}{10}\pi;\ -\frac{103}{10}\pi \right\} \)}{\( \left\{ -\frac{17}{10}\pi;\ -\frac{77}{10}\pi \right\} \)}{\( \left\{ -\frac{37}{10}\pi;\ \frac{23}{10}\pi \right\} \)}{}{}}
\End

\Data\ID{1003023201}
\Updated{2020-12-03 03:12:13}
\Level{A}
\SubArea{Angles, arcs and sectors}
\LANGen{
\Question{
Subtract the angles \( \alpha=1285^{\circ} 35' \), \( \beta= 985^{\circ}59' \)  and round the result to the nearest degree.}
{\( 300^{\circ} \)}{\( 299^{\circ} \)}{\( 301^{\circ} \)}{\( 298^{\circ} \)}{}{}}
\LANGpl{
\Question{
Odejmij kąty   \( \alpha=1285^{\circ} 35' \), \( \beta= 985^{\circ}59' \)  i zaokrągli wynik do pełnych stopni.}
{\( 300^{\circ} \)}{\( 299^{\circ} \)}{\( 301^{\circ} \)}{\( 298^{\circ} \)}{}{}}
\LANGcs{
\Question{
Odečtením úhlů  \( \alpha=1285^{\circ} 35' \), \( \beta= 985^{\circ}59' \) dostaneme úhel zaokrouhlený na celé stupně:}
{\( 300^{\circ} \)}{\( 299^{\circ} \)}{\( 301^{\circ} \)}{\( 298^{\circ} \)}{}{}}
\LANGsk{
\Question{
Odčítaním uhlov   \( \alpha=1285^{\circ} 35' \), \( \beta= 985^{\circ}59' \) dostaneme uhol zaokrúhlený na celé stupne:}
{\( 300^{\circ} \)}{\( 299^{\circ} \)}{\( 301^{\circ} \)}{\( 298^{\circ} \)}{}{}}
\LANGes{
\Question{
Resta los ángulos \( \alpha=1285^{\circ} 35' \), \( \beta= 985^{\circ}59' \)  y aproxima el resultado al grado más cercano.}
{\( 300^{\circ} \)}{\( 299^{\circ} \)}{\( 301^{\circ} \)}{\( 298^{\circ} \)}{}{}}
\End

\Data\ID{1003023202}
\Updated{2020-12-10 10:12:07}
\Level{A}
\SubArea{Angles, arcs and sectors}
\LANGen{
\Question{
Let \( \alpha= 2.7\,\mathrm{rad} \) and \( \beta =54^{\circ} \). The radian measure of the angle \( \alpha-\beta \) (rounded to two decimal places) is:}
{\( 1.76\,\mathrm{rad} \)}{\( -1.76\,\mathrm{rad} \)}{\( 3.65\,\mathrm{rad} \)}{\( -3.65\,\mathrm{rad} \)}{}{}}
\LANGpl{
\Question{
Podano \( \alpha= 2{,}7\,\mathrm{rad} \) i \( \beta =54^{\circ} \). Miara kąta w radianach  \( \alpha-\beta \) (zaokrągli do dwóch miejsc po przecinku) jest równa :}
{\( 1{,}76\,\mathrm{rad} \)}{\( -1{,}76\,\mathrm{rad} \)}{\( 3{,}65\,\mathrm{rad} \)}{\( -3{,}65\,\mathrm{rad} \)}{}{}}
\LANGcs{
\Question{
Jestliže \( \alpha= 2{,}7\,\mathrm{rad} \) a \( \beta =54^{\circ} \), pak velikost úhlu \( \alpha-\beta \) v radiánech (zaokrouhleno na dvě desetinná místa) je:}
{\( 1{,}76\,\mathrm{rad} \)}{\( -1{,}76\,\mathrm{rad} \)}{\( 3{,}65\,\mathrm{rad} \)}{\( -3{,}65\,\mathrm{rad} \)}{}{}}
\LANGsk{
\Question{
Ak  \( \alpha= 2{,}7\,\mathrm{rad} \) and \( \beta =54^{\circ} \), tak veľkosť uhla  \( \alpha-\beta \) v radiánoch (zaokrúhlené na dve desatinné miesta) je:}
{\( 1{,}76\,\mathrm{rad} \)}{\( -1{,}76\,\mathrm{rad} \)}{\( 3{,}65\,\mathrm{rad} \)}{\( -3{,}65\,\mathrm{rad} \)}{}{}}
\LANGes{
\Question{
Sea \( \alpha= 2.7\,\mathrm{rad} \) y \( \beta =54^{\circ} \). La medida del ángulo \( \alpha-\beta \) en radianes (aproximada a dos cifras decimales) es:}
{\( 1.76\,\mathrm{rad} \)}{\( -1.76\,\mathrm{rad} \)}{\( 3.65\,\mathrm{rad} \)}{\( -3.65\,\mathrm{rad} \)}{}{}}
\End

\Data\ID{1003023203}
\Updated{2020-12-10 10:12:31}
\Level{A}
\SubArea{Angles, arcs and sectors}
\LANGen{
\Question{
Let \( \alpha=4.25\,\mathrm{rad} \) and \( \beta = 135^{\circ}15' \). The degree measure of the angle \(\alpha-\beta\) (rounded to the nearest minute) is:}
{\( 108^{\circ}15' \)}{\( 135^{\circ}18' \)}{\( -135^{\circ}18' \)}{\( 405^{\circ}48' \)}{}{}}
\LANGpl{
\Question{
Podano \( \alpha=4{,}25\,\mathrm{rad} \) i  \( \beta = 135^{\circ}15' \). Miara kąta w stopniach  \(\alpha-\beta\) (zaokrągli do pełnych minut ) jest równa:}
{\( 108^{\circ}15' \)}{\( 135^{\circ}18' \)}{\( -135^{\circ}18' \)}{\( 405^{\circ}48' \)}{}{}}
\LANGcs{
\Question{
Jestliže \( \alpha=4{,}25\,\mathrm{rad} \) a \( \beta = 135^{\circ}15' \), pak velikost úhlu \(\alpha-\beta\) v stupních (zaokrouhleno na minuty) je:}
{\( 108^{\circ}15' \)}{\( 135^{\circ}18' \)}{\( -135^{\circ}18' \)}{\( 405^{\circ}48' \)}{}{}}
\LANGsk{
\Question{
Ak   \( \alpha=4{,}25\,\mathrm{rad} \) a \( \beta = 135^{\circ}15' \), tak veľkosť uhla  \(\alpha-\beta\)  v stupňoch (zaokrúhlené na minúty) je:}
{\( 108^{\circ}15' \)}{\( 135^{\circ}18' \)}{\( -135^{\circ}18' \)}{\( 405^{\circ}48' \)}{}{}}
\LANGes{
\Question{
Sea \( \alpha=4.25\,\mathrm{rad} \) y \( \beta = 135^{\circ}15' \). La medida del ángulo \(\alpha-\beta\) en grados (aproximada al minuto más cercano) es:}
{\( 108^{\circ}15' \)}{\( 135^{\circ}18' \)}{\( -135^{\circ}18' \)}{\( 405^{\circ}48' \)}{}{}}
\End

\Data\ID{1003023204}
\Updated{2020-12-10 10:12:19}
\Level{A}
\SubArea{Angles, arcs and sectors}
\LANGen{
\Question{
Let \( \alpha=200^{\circ}15' \) and  \( \beta=  \frac{\pi}5 \).  The degree measure of the angle \( \alpha+\beta \) is:}
{\( 236^{\circ}15' \)}{\( -236^{\circ}15' \)}{\( 201^{\circ}15' \)}{\( -201^{\circ}15' \)}{}{}}
\LANGpl{
\Question{
Podano \( \alpha=200^{\circ}15' \) i  \( \beta=  \frac{\pi}5 \).  Miara kąta w stopniach \( \alpha+\beta \) jest równa:}
{\( 236^{\circ}15' \)}{\( -236^{\circ}15' \)}{\( 201^{\circ}15' \)}{\( -201^{\circ}15' \)}{}{}}
\LANGcs{
\Question{
Pokud platí, že \( \alpha=200^{\circ}15' \) a \( \beta=  \frac{\pi}5 \), potom velikost úhlu \( \alpha+\beta \) v stupních je:}
{\( 236^{\circ}15' \)}{\( -236^{\circ}15' \)}{\( 201^{\circ}15' \)}{\( -201^{\circ}15' \)}{}{}}
\LANGsk{
\Question{
Nech   \( \alpha=200^{\circ}15' \) a  \( \beta=  \frac{\pi}5 \), potom veľkosť uhla    \( \alpha+\beta \) v stupňoch je:}
{\( 236^{\circ}15' \)}{\( -236^{\circ}15' \)}{\( 201^{\circ}15' \)}{\( -201^{\circ}15' \)}{}{}}
\LANGes{
\Question{
Sea \( \alpha=200^{\circ}15' \) y \( \beta=  \frac{\pi}5 \). La medida del ángulo \( \alpha+\beta \) en grados es:}
{\( 236^{\circ}15' \)}{\( -236^{\circ}15' \)}{\( 201^{\circ}15' \)}{\( -201^{\circ}15' \)}{}{}}
\End

\Data\ID{1003023205}
\Updated{2020-12-03 03:12:14}
\Level{A}
\SubArea{Angles, arcs and sectors}
\LANGen{
\Question{
Express the degree measure of an angle if its radian measure is \( \frac{5\pi}{12} \).}
{\( 75^{\circ} \)}{\( 105^{\circ} \)}{\( 225^{\circ} \)}{\( 285^{\circ} \)}{}{}}
\LANGpl{
\Question{
Wyraź miarę kąta w stopniach, jeśli jego  miara podana w radianach jest równa \( \frac{5\pi}{12} \).}
{\( 75^{\circ} \)}{\( 105^{\circ} \)}{\( 225^{\circ} \)}{\( 285^{\circ} \)}{}{}}
\LANGcs{
\Question{
Vyjádřete velikost úhlu ve stupních, jestliže jeho velikost v radiánech je \( \frac{5\pi}{12} \).}
{\( 75^{\circ} \)}{\( 105^{\circ} \)}{\( 225^{\circ} \)}{\( 285^{\circ} \)}{}{}}
\LANGsk{
\Question{
Vyjadrite veľkosť uhla v stupňoch, ak jeho veľkosť v radiánoch je   \( \frac{5\pi}{12} \).}
{\( 75^{\circ} \)}{\( 105^{\circ} \)}{\( 225^{\circ} \)}{\( 285^{\circ} \)}{}{}}
\LANGes{
\Question{
Expresa la medida del un ángulo en grados si su medida en radianes es \( \frac{5\pi}{12} \).}
{\( 75^{\circ} \)}{\( 105^{\circ} \)}{\( 225^{\circ} \)}{\( 285^{\circ} \)}{}{}}
\End

\Data\ID{1003023301}
\Updated{2021-06-21 09:06:30}
\Level{A}
\SubArea{Angles, arcs and sectors}
\LANGen{
\Question{
Which value is not the measure of an angle that is coterminal with an angle of radian measure \( \frac23\pi \)? (Two angles are coterminal if they are drawn in the standard position and both have their terminal sides in the same location.)}
{\( -\frac13\pi \)}{\( -\frac43\pi \)}{\( \frac83\pi \)}{\( -\frac{10}3\pi \)}{}{}}
\LANGpl{
\Question{
Która wartość nie jest równa kątowi \( \frac23\pi \).}
{\( -\frac13\pi \)}{\( -\frac43\pi \)}{\( \frac83\pi \)}{\( -\frac{10}3\pi \)}{}{}}
\LANGcs{
\Question{
Orientovaný úhel má základní velikost \( \frac23\pi \). Které z následujících čísel nepředstavuje velikost tohoto orientovaného úhlu?}
{\( -\frac13\pi \)}{\( -\frac43\pi \)}{\( \frac83\pi \)}{\( -\frac{10}3\pi \)}{}{}}
\LANGsk{
\Question{
Orientovaný uhol má základnú veľkosť   \( \frac23\pi \). Ktoré z nasledujúcich čísel nepredstavuje veľkosť tohto orientovaného uhla?}
{\( -\frac13\pi \)}{\( -\frac43\pi \)}{\( \frac83\pi \)}{\( -\frac{10}3\pi \)}{}{}}
\LANGes{
\Question{
Elige la medida del ángulo coterminal al ángulo \( \frac23\pi \) radianes.  (Dos ángulos son coterminales si dibujándolos en la posición estándar tienen el lado terminal en el mismo sitio)}
{\( -\frac13\pi \)}{\( -\frac43\pi \)}{\( \frac83\pi \)}{\( -\frac{10}3\pi \)}{}{}}
\End

\Data\ID{1003023302}
\Updated{2021-06-21 09:06:49}
\Level{A}
\SubArea{Angles, arcs and sectors}
\LANGen{
\Question{
Select the measure of an angle that is coterminal to an angle of \( 40^{\circ} \). (Two angles are coterminal if they are drawn in the standard position and both have their terminal sides in the same location.)}
{\( 400^{\circ} \)}{\( 440^{\circ} \)}{\( 360^{\circ} \)}{\( 800^{\circ} \)}{}{}}
\LANGpl{
\Question{
Wybierz miarę kąta skierowanego równemu  kątowi \( 40^{\circ} \)?}
{\( 400^{\circ} \)}{\( 440^{\circ} \)}{\( 360^{\circ} \)}{\( 800^{\circ} \)}{}{}}
\LANGcs{
\Question{
Orientovaný úhel má základní velikost \( 40^{\circ} \). Z následujících možností vyberte jinou velikost stejného orientovaného úhlu.}
{\( 400^{\circ} \)}{\( 440^{\circ} \)}{\( 360^{\circ} \)}{\( 800^{\circ} \)}{}{}}
\LANGsk{
\Question{
Orientovaný uhol má základnú veľkosť   \( 40^{\circ} \). Z nasledujúcej ponuky vyberte inú veľkosť tohto orientovaného uhla.}
{\( 400^{\circ} \)}{\( 440^{\circ} \)}{\( 360^{\circ} \)}{\( 800^{\circ} \)}{}{}}
\LANGes{
\Question{
Elige la medida del ángulo coterminal al ángulo \( 40^{\circ} \). (Dos ángulos son coterminales si dibujándolos en la posición estándar tienen el lado terminal en el mismo sitio)}
{\( 400^{\circ} \)}{\( 440^{\circ} \)}{\( 360^{\circ} \)}{\( 800^{\circ} \)}{}{}}
\End

\Data\ID{1003023303}
\Updated{2020-12-10 03:12:24}
\Level{A}
\SubArea{Angles, arcs and sectors}
\LANGen{
\Question{
Which value is the measure of an angle that is coterminal with an angle of radian measure \( \frac{23}3\pi \)? (Two angles are coterminal if they are drawn in the standard position and both have their terminal sides in the same location.)}
{\( \frac53\pi \)}{\( \frac13\pi \)}{\( -\frac53\pi \)}{\( -\frac43\pi \)}{}{}}
\LANGpl{
\Question{
Wybierz miarę kąta skierowanego równemu kątowi \( \frac{23}3\pi \).}
{\( \frac53\pi \)}{\( \frac13\pi \)}{\( -\frac53\pi \)}{\( -\frac43\pi \)}{}{}}
\LANGcs{
\Question{
Jedna velikost orientovaného úhlu je \( \frac{23}3\pi \). Jaká je jeho základní velikost?}
{\( \frac53\pi \)}{\( \frac13\pi \)}{\( -\frac53\pi \)}{\( -\frac43\pi \)}{}{}}
\LANGsk{
\Question{
Jedna veľkosť orientovaného uhla je \( \frac{23}3\pi \). Aká je jeho základná veľkosť ?}
{\( \frac53\pi \)}{\( \frac13\pi \)}{\( -\frac53\pi \)}{\( -\frac43\pi \)}{}{}}
\LANGes{
\Question{
¿Qué valor es la medida del ángulo que es coterminal con el ángulo de \( \frac{23}3\pi \) radianes? (Dos ángulos son coterminales si dibujándolos en la posición estándar tienen los lados terminales en el mismo sitio.)}
{\( \frac53\pi \)}{\( \frac13\pi \)}{\( -\frac53\pi \)}{\( -\frac43\pi \)}{}{}}
\End

\Data\ID{1003023304}
\Updated{2020-12-10 02:12:41}
\Level{A}
\SubArea{Angles, arcs and sectors}
\LANGen{
\Question{
Which of the given degree values is of the angle coterminal to the angle of \( -2000^{\circ} \)? (Two angles are coterminal if they are drawn in the standard position and both have their terminal sides in the same location.)}
{\( 160^{\circ} \)}{\(-160^{\circ} \)}{\( -20^{\circ} \)}{\( 200^{\circ} \)}{}{}}
\LANGpl{
\Question{
Wybierz miarę kąta skierowanego równemu kątowi \( -2000^{\circ} \).}
{\( 160^{\circ} \)}{\(-160^{\circ} \)}{\( -20^{\circ} \)}{\( 200^{\circ} \)}{}{}}
\LANGcs{
\Question{
Jedna velikost orientovaného úhlu je  \( -2000^{\circ} \). Jaká je jeho základní velikost?}
{\( 160^{\circ} \)}{\(-160^{\circ} \)}{\( -20^{\circ} \)}{\( 200^{\circ} \)}{}{}}
\LANGsk{
\Question{
Jedna veľkosť orientovaného uhla je \( -2000^{\circ} \). Aká je jeho základná veľkosť ?}
{\( 160^{\circ} \)}{\(-160^{\circ} \)}{\( -20^{\circ} \)}{\( 200^{\circ} \)}{}{}}
\LANGes{
\Question{
¿Qué valor es la medida del ángulo que es coterminal con el ángulo de  \( -2000^{\circ} \)? (Dos ángulos son coterminales si dibujándolos en la posición estándar tienen los lados terminales en el mismo sitio.)}
{\( 160^{\circ} \)}{\(-160^{\circ} \)}{\( -20^{\circ} \)}{\( 200^{\circ} \)}{}{}}
\End

\Data\ID{1003023305}
\Updated{2020-12-10 02:12:56}
\Level{A}
\SubArea{Angles, arcs and sectors}
\LANGen{
\Question{
Select the set that does not contain the measures of coterminal angles with an angle of radian measure \( \frac{\pi}3 \). (Two angles are coterminal if they are drawn in the standard position and both have their terminal sides in the same location.)}
{\( \left\{\frac43\pi;\ -\frac{10}3\pi \right\} \)}{\( \left\{\frac73\pi;\ -\frac53\pi \right\} \)}{\( \left\{\frac{13}3\pi;\ \frac{61}3\pi \right\} \)}{\( \left\{\frac{19}3\pi;\ \frac{25}3\pi \right\} \)}{}{}}
\LANGpl{
\Question{
Do którego zbioru nie należy kąt \( \frac{\pi}3 \).}
{\( \left\{\frac43\pi;\ -\frac{10}3\pi \right\} \)}{\( \left\{\frac73\pi;\ -\frac53\pi \right\} \)}{\( \left\{\frac{13}3\pi;\ \frac{61}3\pi \right\} \)}{\( \left\{\frac{19}3\pi;\ \frac{25}3\pi \right\} \)}{}{}}
\LANGcs{
\Question{
Základní velikost orientovaného úhlu je \( \frac{\pi}3 \).  Z následujících možností vyberte množinu, která neobsahuje velikost tohoto orientovaného úhlu.}
{\( \left\{\frac43\pi;\ -\frac{10}3\pi \right\} \)}{\( \left\{\frac73\pi;\ -\frac53\pi \right\} \)}{\( \left\{\frac{13}3\pi;\ \frac{61}3\pi \right\} \)}{\( \left\{\frac{19}3\pi;\ \frac{25}3\pi \right\} \)}{}{}}
\LANGsk{
\Question{
Základná veľkosť orientovaného uhla je \( \frac{\pi}3 \).  Z nasledujúcej ponuky vyberte tú množinu, ktorá neobsahuje veľkosti tohto orientovaného uhla.}
{\( \left\{\frac43\pi;\ -\frac{10}3\pi \right\} \)}{\( \left\{\frac73\pi;\ -\frac53\pi \right\} \)}{\( \left\{\frac{13}3\pi;\ \frac{61}3\pi \right\} \)}{\( \left\{\frac{19}3\pi;\ \frac{25}3\pi \right\} \)}{}{}}
\LANGes{
\Question{
Elije el conjunto que no contiene medidas de ángulos coterminales con el ángulo de \( \frac{\pi}3 \) radianes.(Dos ángulos son coterminales si dibujándolos en la posición estándar tienen los lados terminales en el mismo sitio.)}
{\( \left\{\frac43\pi;\ -\frac{10}3\pi \right\} \)}{\( \left\{\frac73\pi;\ -\frac53\pi \right\} \)}{\( \left\{\frac{13}3\pi;\ \frac{61}3\pi \right\} \)}{\( \left\{\frac{19}3\pi;\ \frac{25}3\pi \right\} \)}{}{}}
\End

\Data\ID{1003023310}
\Updated{2020-12-10 03:12:53}
\Level{A}
\SubArea{Angles, arcs and sectors}
\LANGen{
\Question{
Which of the given radian measures is of the angle coterminal to the angle of \( -50\pi \)? (Two angles are coterminal if they are drawn in the standard position and both have their terminal sides in the same location.)}
{\( 0 \)}{\( \pi \)}{\( 3\pi \)}{\( -\pi \)}{}{}}
\LANGpl{
\Question{
Podaj miarę kąta skierowanego równego kątowi  \( -50\pi \)?}
{\( 0 \)}{\( \pi \)}{\( 3\pi \)}{\( -\pi \)}{}{}}
\LANGcs{
\Question{
Základní velikost orientovaného úhlu \( -50\pi \) je:}
{\( 0 \)}{\( \pi \)}{\( 3\pi \)}{\( -\pi \)}{}{}}
\LANGsk{
\Question{
Základná veľkosť orientovaného uhla \( -50\pi \) je:}
{\( 0 \)}{\( \pi \)}{\( 3\pi \)}{\( -\pi \)}{}{}}
\LANGes{
\Question{
¿Cuál de las medidas en radianes es la del ángulo coterminal a \( -50\pi \)? (Dos ángulos son coterminales si dibujándolos en la posición estándar tienen los lados terminales en el mismo sitio.)}
{\( 0 \)}{\( \pi \)}{\( 3\pi \)}{\( -\pi \)}{}{}}
\End

\Data\ID{1003023311}
\Updated{2020-12-10 03:12:11}
\Level{A}
\SubArea{Angles, arcs and sectors}
\LANGen{
\Question{
Which of the given degree values is not of the angle coterminal to the angle of \( 60^{\circ} \)? (Two angles are coterminal if they are drawn in the standard position and both have their terminal sides in the same location.)}
{\( 300^{\circ} \)}{\( -300^{\circ} \)}{\( -660^{\circ} \)}{\( 420^{\circ} \)}{}{}}
\LANGpl{
\Question{
Które z poniższych wartości nie są równe kątowi skierowanemu  równemu  \( 60^{\circ} \)?}
{\( 300^{\circ} \)}{\( -300^{\circ} \)}{\( -660^{\circ} \)}{\( 420^{\circ} \)}{}{}}
\LANGcs{
\Question{
Který z daných orientovaných úhlů nemá základní velikost \( 60^{\circ} \)?}
{\( 300^{\circ} \)}{\( -300^{\circ} \)}{\( -660^{\circ} \)}{\( 420^{\circ} \)}{}{}}
\LANGsk{
\Question{
Ktorý z daných orientovaných uhlov nemá základnú veľkosť  \( 60^{\circ} \)?}
{\( 300^{\circ} \)}{\( -300^{\circ} \)}{\( -660^{\circ} \)}{\( 420^{\circ} \)}{}{}}
\LANGes{
\Question{
¿Cuál de las medidas en grados no es el valor del ángulo coterminal al ángulo \( 60^{\circ} \)?(Dos ángulos son coterminales si dibujándolos en la posición estándar tienen los lados terminales en el mismo sitio.)}
{\( 300^{\circ} \)}{\( -300^{\circ} \)}{\( -660^{\circ} \)}{\( 420^{\circ} \)}{}{}}
\End

\Data\ID{1003023401}
\Updated{2020-12-10 03:12:02}
\Level{A}
\SubArea{Angles, arcs and sectors}
\LANGen{
\Question{
An angle of radian measure \( 1 \) is in standard position. In which quadrant lies its terminal side?}
{I.}{II.}{III.}{IV.}{}{}}
\LANGpl{
\Question{
W której ćwiartce układu współrzędnych leży ramię końcowe kąta, którego miarą jest \( 1 \) radian?}
{I.}{II.}{III.}{IV.}{}{}}
\LANGcs{
\Question{
Ve kterém kvadrantu leží koncové rameno orientovaného úhlu \( 1 \)    rad?}
{I.}{II.}{III.}{IV.}{}{}}
\LANGsk{
\Question{
V ktorom kvadrante leží koncové rameno orientovaného uhla  \( 1 \)    rad?}
{I.}{II.}{III.}{IV.}{}{}}
\LANGes{
\Question{
El ángulo que mide \( 1 \) radián es en posición estándar. ¿En cuál de los cuadrantes está su lado terminal?}
{I.}{II.}{III.}{IV.}{}{}}
\End

\Data\ID{1103023106}
\Updated{2021-01-13 12:01:09}
\Level{A}
\SubArea{Angles, arcs and sectors}
\IMG{1103023106.pdf}
\LANGen{
\Question{
Which of the angles \( \alpha \), \( \beta \), \( \gamma \), \( \delta \) takes the same position on the unit circle as the angle \( APB \)?}
{\( \alpha = 123^{\circ} \)}{\( \beta = -100^{\circ} \)}{\( \gamma = -340^{\circ} \)}{\( \delta = -380^{\circ} \)}{}{}}
\LANGpl{
\Question{
Który z kątów \( \alpha \), \( \beta \), \( \gamma \), \( \delta \) przyjmuje tę samą pozycję na okręgu jednostkowym, co kąt  \( APB \)?}
{\( \alpha = 123^{\circ} \)}{\( \beta = -100^{\circ} \)}{\( \gamma = -340^{\circ} \)}{\( \delta = -380^{\circ} \)}{}{}}
\LANGcs{
\Question{
Který z úhlů  \( \alpha \), \( \beta \), \( \gamma \), \( \delta \) se zobrazí na jednotkové kružnici stejně jako úhel \( APB \)?}
{\( \alpha = 123^{\circ} \)}{\( \beta = -100^{\circ} \)}{\( \gamma = -340^{\circ} \)}{\( \delta = -380^{\circ} \)}{}{}}
\LANGsk{
\Question{
Ktorý z uhlov  \( \alpha \), \( \beta \), \( \gamma \), \( \delta \) sa zobrazí na jednotkovej kružnici rovnako ako uhol  \( APB \)?}
{\( \alpha = 123^{\circ} \)}{\( \beta = -100^{\circ} \)}{\( \gamma = -340^{\circ} \)}{\( \delta = -380^{\circ} \)}{}{}}
\LANGes{
\Question{
¿Cuál de los ángulos  \( \alpha \), \( \beta \), \( \gamma \), \( \delta \) aparece en la misma posición en la circunferencia unitaria que el ángulo \( APB \)?}
{\( \alpha = 123^{\circ} \)}{\( \beta = -100^{\circ} \)}{\( \gamma = -340^{\circ} \)}{\( \delta = -380^{\circ} \)}{}{}}
\End

\Data\ID{1103023107}
\Updated{2021-01-13 12:01:15}
\Level{A}
\SubArea{Angles, arcs and sectors}
\IMG{1103023107.pdf}
\LANGen{
\Question{
Which of the angles \( \alpha \), \( \beta \), \( \gamma \), \( \delta \) takes the same position on the unit circle as the angle \( APB \)?}
{\( \gamma = 4.18\,\mathrm{rad} \)}{\( \beta = 7.32\,\mathrm{rad} \)}{\( \alpha = -7.32\,\mathrm{rad} \)}{\( \delta = -4.18\,\mathrm{rad} \)}{}{}}
\LANGpl{
\Question{
Który z kątów \( \alpha \), \( \beta \), \( \gamma \), \( \delta \) przyjmuje tę samą pozycję na okręgu jednostkowym, co kąt \( APB \)?}
{\( \gamma = 4{,}18\,\mathrm{rad} \)}{\( \beta = 7{,}32\,\mathrm{rad} \)}{\( \alpha = -7{,}32\,\mathrm{rad} \)}{\( \delta = -4{,}18\,\mathrm{rad} \)}{}{}}
\LANGcs{
\Question{
Který z úhlů \( \alpha \), \( \beta \), \( \gamma \), \( \delta \) se zobrazí na jednotkové kružnici stejně jako úhel \( APB \)?}
{\( \gamma = 4{,}18\,\mathrm{rad} \)}{\( \beta = 7{,}32\,\mathrm{rad} \)}{\( \alpha = -7{,}32\,\mathrm{rad} \)}{\( \delta = -4{,}18\,\mathrm{rad} \)}{}{}}
\LANGsk{
\Question{
Ktorý z uhlov  \( \alpha \), \( \beta \), \( \gamma \), \( \delta \) sa zobrazí na jednotkovej kružnici rovnako ako uhol \( APB \)?}
{\( \gamma = 4{,}18\,\mathrm{rad} \)}{\( \beta = 7{,}32\,\mathrm{rad} \)}{\( \alpha = -7{,}32\,\mathrm{rad} \)}{\( \delta = -4{,}18\,\mathrm{rad} \)}{}{}}
\LANGes{
\Question{
¿Cuál de los ángulos \( \alpha \), \( \beta \), \( \gamma \), \( \delta \) aparece en la misma posición en la circunferencia unitaria que el ángulo \( APB \)?}
{\( \gamma = 4{,}18\,\mathrm{rad} \)}{\( \beta = 7{,}32\,\mathrm{rad} \)}{\( \alpha = -7{,}32\,\mathrm{rad} \)}{\( \delta = -4{,}18\,\mathrm{rad} \)}{}{}}
\End

\Data\ID{9000033901}
\Updated{2021-01-10 01:01:12}
\Level{A}
\SubArea{Angles, arcs and sectors}
\LANGen{
\Question{
Associate the angle \(\frac{7}
{6}\pi \) 
to the corresponding quadrant.}
{III.}{I.}{II.}{IV.}{}{}}
\LANGpl{
\Question{
Przyporządkuj kąt  \(\frac{7}
{6}\pi \) 
do odpowiedniej ćwiartki układu współrzędnych.}
{III.}{I.}{II.}{IV.}{}{}}
\LANGcs{
\Question{
Přiřaďte úhlu \(\frac{7}
{6}\pi \) 
příslušný kvadrant.}
{III.}{I.}{II.}{IV.}{}{}}
\LANGsk{
\Question{
Priraďte uhlu \(\frac{7}
{6}\pi \) 
príslušný kvadrant.}
{III.}{I.}{II.}{IV.}{}{}}
\LANGes{
\Question{
¿En qué cuadrante está el ángulo \(\frac{7}
{6}\pi \) ?}
{III.}{I.}{II.}{IV.}{}{}}
\End

\Data\ID{9000033902}
\Updated{2021-01-10 01:01:32}
\Level{A}
\SubArea{Angles, arcs and sectors}
\LANGen{
\Question{
Associate the angle \(\frac{7}
{8}\pi \) 
to the corresponding quadrant.}
{II.}{I.}{III.}{IV.}{}{}}
\LANGpl{
\Question{
Przyporządkuj kąt  \(\frac{7}
{8}\pi \) 
do odpowiedniej ćwiartki układu współrzędnych.}
{II.}{I.}{III.}{IV.}{}{}}
\LANGcs{
\Question{
Přiřaďte úhlu \(\frac{7}
{8}\pi \) 
příslušný kvadrant.}
{II.}{I.}{III.}{IV.}{}{}}
\LANGsk{
\Question{
Priraďte uhlu \(\frac{7}
{8}\pi \) 
príslušný kvadrant.}
{II.}{I.}{III.}{IV.}{}{}}
\LANGes{
\Question{
¿En qué cuadrante está el ángulo \(\frac{7}
{8}\pi \) ?}
{II.}{I.}{III.}{IV.}{}{}}
\End

\Data\ID{9000033903}
\Updated{2021-01-13 12:01:29}
\Level{A}
\SubArea{Angles, arcs and sectors}
\LANGen{
\Question{
In the interval \([0;2\pi )\) find the angle
equivalent to the angle \(\frac{21}
 {6} \pi \).}
{\(\frac{3}
{2}\pi \)}{\(\frac{\pi }{2}\)}{\(\frac{\pi }
{3}\)}{\(\frac{2}
{3}\pi \)}{}{}}
\LANGpl{
\Question{
Na przedziale  \([0;2\pi )\) wskaż kąt odpowiadający do kąta  \(\frac{21}
 {6} \pi \).}
{\(\frac{3}
{2}\pi \)}{\(\frac{\pi }{2}\)}{\(\frac{\pi }
{3}\)}{\(\frac{2}
{3}\pi \)}{}{}}
\LANGcs{
\Question{
Základní velikost úhlu \(\frac{21}
 {6} \pi \) 
je:}
{\(\frac{3}
{2}\pi \)}{\(\frac{\pi }{2}\)}{\(\frac{\pi }
{3}\)}{\(\frac{2}
{3}\pi \)}{}{}}
\LANGsk{
\Question{
Základná veľkosť uhla \(\frac{21}
 {6} \pi \) 
je:}
{\(\frac{3}
{2}\pi \)}{\(\frac{\pi }{2}\)}{\(\frac{\pi }
{3}\)}{\(\frac{2}
{3}\pi \)}{}{}}
\LANGes{
\Question{
La medida canónica del ángulo\(\frac{21}
 {6} \pi \) 
es:}
{\(\frac{3}
{2}\pi \)}{\(\frac{\pi }{2}\)}{\(\frac{\pi }
{3}\)}{\(\frac{2}
{3}\pi \)}{}{}}
\End

\Data\ID{9000033904}
\Updated{2021-01-04 07:01:37}
\Level{A}
\SubArea{Angles, arcs and sectors}
\LANGen{
\Question{
In the interval \([0;2\pi )\) find the angle
equivalent to the angle \(-\frac{17} 
 {3} \pi \).}
{\(\frac{\pi }{3}\)}{\(\frac{2}
{3}\pi \)}{\(\frac{4}
{3}\pi \)}{\(\frac{5}
{3}\pi \)}{}{}}
\LANGpl{
\Question{
Na przedziale \([0;2\pi )\) wskaż kąt odpowiadający do kąta \(-\frac{17} 
 {3} \pi \).}
{\(\frac{\pi }{3}\)}{\(\frac{2}
{3}\pi \)}{\(\frac{4}
{3}\pi \)}{\(\frac{5}
{3}\pi \)}{}{}}
\LANGcs{
\Question{
Základní velikost úhlu \(-\frac{17} 
 {3} \pi \) 
je:}
{\(\frac{\pi }{3}\)}{\(\frac{2}
{3}\pi \)}{\(\frac{4}
{3}\pi \)}{\(\frac{5}
{3}\pi \)}{}{}}
\LANGsk{
\Question{
Základná veľkosť uhla \(-\frac{17} 
 {3} \pi \) 
je:}
{\(\frac{\pi }{3}\)}{\(\frac{2}
{3}\pi \)}{\(\frac{4}
{3}\pi \)}{\(\frac{5}
{3}\pi \)}{}{}}
\LANGes{
\Question{
La posición canónica del ángulo \(-\frac{17} 
 {3} \pi \) 
es:}
{\(\frac{\pi }{3}\)}{\(\frac{2}
{3}\pi \)}{\(\frac{4}
{3}\pi \)}{\(\frac{5}
{3}\pi \)}{}{}}
\End

\Data\ID{9000033905}
\Updated{2021-01-04 07:01:56}
\Level{A}
\SubArea{Angles, arcs and sectors}
\LANGen{
\Question{
In the interval \([0^{\circ };360^{\circ })\) find the angle
equivalent to the angle \(- 428^{\circ }\).}
{\(292^{\circ }\)}{\(192^{\circ }\)}{\(68^{\circ }\)}{\(168^{\circ }\)}{}{}}
\LANGpl{
\Question{
Na przedziale  \([0^{\circ };360^{\circ })\) wskaż kąt odpowiadający do kąta \(- 428^{\circ }\).}
{\(292^{\circ }\)}{\(192^{\circ }\)}{\(68^{\circ }\)}{\(168^{\circ }\)}{}{}}
\LANGcs{
\Question{
Základní velikost úhlu \(- 428^{\circ }\) 
je:}
{\(292^{\circ }\)}{\(192^{\circ }\)}{\(68^{\circ }\)}{\(168^{\circ }\)}{}{}}
\LANGsk{
\Question{
Základná veľkosť uhla \(- 428^{\circ }\) 
je:}
{\(292^{\circ }\)}{\(192^{\circ }\)}{\(68^{\circ }\)}{\(168^{\circ }\)}{}{}}
\LANGes{
\Question{
La posición canónica del ángulo \(- 428^{\circ }\) es:}
{\(292^{\circ }\)}{\(192^{\circ }\)}{\(68^{\circ }\)}{\(168^{\circ }\)}{}{}}
\End

\Data\ID{9000033906}
\Updated{2021-01-04 07:01:12}
\Level{A}
\SubArea{Angles, arcs and sectors}
\LANGen{
\Question{
In the interval \([0^{\circ };360^{\circ })\) find the angle
equivalent to the angle \(1\: 000^{\circ }\).}
{\(280^{\circ }\)}{\(180^{\circ }\)}{\(240^{\circ }\)}{\(300^{\circ }\)}{}{}}
\LANGpl{
\Question{
Na przedziale \([0^{\circ };360^{\circ })\) wskaż kąt odpowiadający do kąta \(1\: 000^{\circ }\).}
{\(280^{\circ }\)}{\(180^{\circ }\)}{\(240^{\circ }\)}{\(300^{\circ }\)}{}{}}
\LANGcs{
\Question{
Základní velikost úhlu \(1\: 000^{\circ }\) 
je:}
{\(280^{\circ }\)}{\(180^{\circ }\)}{\(240^{\circ }\)}{\(300^{\circ }\)}{}{}}
\LANGsk{
\Question{
Základná veľkosť uhla \(1\: 000^{\circ }\) 
je:}
{\(280^{\circ }\)}{\(180^{\circ }\)}{\(240^{\circ }\)}{\(300^{\circ }\)}{}{}}
\LANGes{
\Question{
La posición canónica del ángulo \(1\: 000^{\circ }\) 
es:}
{\(280^{\circ }\)}{\(180^{\circ }\)}{\(240^{\circ }\)}{\(300^{\circ }\)}{}{}}
\End

\Data\ID{9000033907}
\Updated{2021-01-04 07:01:30}
\Level{A}
\SubArea{Angles, arcs and sectors}
\LANGen{
\Question{
Convert \(\frac{6}
{5}\pi \) 
radians to degrees.}
{\(216^{\circ }\)}{\(432^{\circ }\)}{\(116^{\circ }\)}{\(378^{\circ }\)}{}{}}
\LANGpl{
\Question{
Podaj miarę kąta  \(\frac{6}
{5}\pi \) 
w stopniach.}
{\(216^{\circ }\)}{\(432^{\circ }\)}{\(116^{\circ }\)}{\(378^{\circ }\)}{}{}}
\LANGcs{
\Question{
Velikost úhlu \(\frac{6}
{5}\pi \) 
v míře stupňové je:}
{\(216^{\circ }\)}{\(432^{\circ }\)}{\(116^{\circ }\)}{\(378^{\circ }\)}{}{}}
\LANGsk{
\Question{
Veľkosť uhla \(\frac{6}
{5}\pi \) 
v miere stupňovej je:}
{\(216^{\circ }\)}{\(432^{\circ }\)}{\(116^{\circ }\)}{\(378^{\circ }\)}{}{}}
\LANGes{
\Question{
El tamaño del ángulo \(\frac{6}
{5}\pi \) 
en grados es:}
{\(216^{\circ }\)}{\(432^{\circ }\)}{\(116^{\circ }\)}{\(378^{\circ }\)}{}{}}
\End

\Data\ID{9000033908}
\Updated{2021-01-04 07:01:32}
\Level{A}
\SubArea{Angles, arcs and sectors}
\LANGen{
\Question{
Convert \(\frac{8}
{3}\pi \) 
radians to degrees.}
{\(480^{\circ }\)}{\(240^{\circ }\)}{\(300^{\circ }\)}{\(330^{\circ }\)}{}{}}
\LANGpl{
\Question{
Podaj miarę kąta \(\frac{8}
{3}\pi \) 
w stopniach.}
{\(480^{\circ }\)}{\(240^{\circ }\)}{\(300^{\circ }\)}{\(330^{\circ }\)}{}{}}
\LANGcs{
\Question{
Velikost úhlu \(\frac{8}
{3}\pi \) 
v míře stupňové je:}
{\(480^{\circ }\)}{\(240^{\circ }\)}{\(300^{\circ }\)}{\(330^{\circ }\)}{}{}}
\LANGsk{
\Question{
Veľkosť uhla \(\frac{8}
{3}\pi \) 
v miere stupňovej je:}
{\(480^{\circ }\)}{\(240^{\circ }\)}{\(300^{\circ }\)}{\(330^{\circ }\)}{}{}}
\LANGes{
\Question{
El tamaño del ángulo \(\frac{8}
{3}\pi \) 
en grados es:}
{\(480^{\circ }\)}{\(240^{\circ }\)}{\(300^{\circ }\)}{\(330^{\circ }\)}{}{}}
\End

\Data\ID{9000033909}
\Updated{2021-01-04 07:01:48}
\Level{A}
\SubArea{Angles, arcs and sectors}
\LANGen{
\Question{
Convert the angle \(240^{\circ }\) 
to radians.}
{\(\frac{4}
{3}\pi \)}{\(\frac{8}
{3}\pi \)}{\(\frac{10}
 {6} \pi \)}{\(\frac{5}
{3}\pi \)}{}{}}
\LANGpl{
\Question{
Podaj miarę kąta  \(240^{\circ }\) 
w radianach.}
{\(\frac{4}
{3}\pi \)}{\(\frac{8}
{3}\pi \)}{\(\frac{10}
 {6} \pi \)}{\(\frac{5}
{3}\pi \)}{}{}}
\LANGcs{
\Question{
Velikost úhlu \(240^{\circ }\) 
v míře obloukové je:}
{\(\frac{4}
{3}\pi \)}{\(\frac{8}
{3}\pi \)}{\(\frac{10}
 {6} \pi \)}{\(\frac{5}
{3}\pi \)}{}{}}
\LANGsk{
\Question{
Veľkosť uhla \(240^{\circ }\) 
v miere oblúkovej je:}
{\(\frac{4}
{3}\pi \)}{\(\frac{8}
{3}\pi \)}{\(\frac{10}
 {6} \pi \)}{\(\frac{5}
{3}\pi \)}{}{}}
\LANGes{
\Question{
El tamaño del ángulo \(240^{\circ }\) 
en radianes es:}
{\(\frac{4}
{3}\pi \)}{\(\frac{8}
{3}\pi \)}{\(\frac{10}
 {6} \pi \)}{\(\frac{5}
{3}\pi \)}{}{}}
\End

\Data\ID{9000033910}
\Updated{2021-01-04 07:01:12}
\Level{A}
\SubArea{Angles, arcs and sectors}
\LANGen{
\Question{
Convert the angle \(292.5^{\circ }\) 
to radians.}
{\(\frac{13}
 {8} \pi \)}{\(\frac{11}
 {4} \pi \)}{\(\frac{15}
 {8} \pi \)}{\(\frac{13}
 {4} \pi \)}{}{}}
\LANGpl{
\Question{
Podaj miarę kąta  \(292.5^{\circ }\) 
w radianach.}
{\(\frac{13}
 {8} \pi \)}{\(\frac{11}
 {4} \pi \)}{\(\frac{15}
 {8} \pi \)}{\(\frac{13}
 {4} \pi \)}{}{}}
\LANGcs{
\Question{
Velikost úhlu \(292{,}5^{\circ }\) 
v míře obloukové je:}
{\(\frac{13}
 {8} \pi \)}{\(\frac{11}
 {4} \pi \)}{\(\frac{15}
 {8} \pi \)}{\(\frac{13}
 {4} \pi \)}{}{}}
\LANGsk{
\Question{
Veľkosť uhla \(292{,}5^{\circ }\) 
v miere oblúkovej je:}
{\(\frac{13}
 {8} \pi \)}{\(\frac{11}
 {4} \pi \)}{\(\frac{15}
 {8} \pi \)}{\(\frac{13}
 {4} \pi \)}{}{}}
\LANGes{
\Question{
Convierte el ángulo \(292{,}5^{\circ }\) a radianes.}
{\(\frac{13}
 {8} \pi \)}{\(\frac{11}
 {4} \pi \)}{\(\frac{15}
 {8} \pi \)}{\(\frac{13}
 {4} \pi \)}{}{}}
\End

\Data\ID{1003023206}
\Updated{2021-06-21 09:06:00}
\Level{B}
\SubArea{Angles, arcs and sectors}
\LANGen{
\Question{
In a  driving school we learn  to hold the steering wheel in the position of „a quarter to three“. Determine the degree measure of the angle between the hour hand and the minute hand which corresponds to the mentioned time.}
{\( 172.5^{\circ} \)}{\( 180^{\circ} \)}{\( 165^{\circ} \)}{\( 157.5^{\circ} \)}{}{}}
\LANGpl{
\Question{
W szkole nauki jazdy uczymy się trzymać kierownicę w pozycji  kwadrans do trzeciej. Określ miarę stopnia kąta między wskazówką godzinową a wskazówką minutową, która odpowiada wspomnianej godzinie.}
{\( 172{,}5^{\circ} \)}{\( 180^{\circ} \)}{\( 165^{\circ} \)}{\( 157{,}5^{\circ} \)}{}{}}
\LANGcs{
\Question{
V autoškole se učíme držet volant v poloze „tři čtvrtě na tři“. Vypočítejte velikost úhlu mezi hodinovou a minutovou ručičkou, které ukazují právě tento čas.}
{\( 172{,}5^{\circ} \)}{\( 180^{\circ} \)}{\( 165^{\circ} \)}{\( 157{,}5^{\circ} \)}{}{}}
\LANGsk{
\Question{
V autoškole sa učíme, že volant treba držať v polohe „trištvrte na tri“. Určte v stupňoch veľkosť uhla medzi hodinovou a minútovou ručičkou, ktorý odpovedá tomuto času.}
{\( 172{,}5^{\circ} \)}{\( 180^{\circ} \)}{\( 165^{\circ} \)}{\( 157{,}5^{\circ} \)}{}{}}
\LANGes{
\Question{
En la autoescuela nos enseñan a sujetar el volante en la posición "tres menos cuatro“. Calcula el ángulo entre la manilla pequeña y la grande en este tiempo.}
{\( 172{,}5^{\circ} \)}{\( 180^{\circ} \)}{\( 165^{\circ} \)}{\( 157{,}5^{\circ} \)}{}{}}
\End

\Data\ID{1003023207}
\Updated{2021-06-21 09:06:00}
\Level{B}
\SubArea{Angles, arcs and sectors}
\LANGen{
\Question{
In the past we were taught  in a  driving school that the steering wheel should be held in the position of „ten to two“. Give the degree measure of the angle between the hour hand and the minute hand which corresponds to this time.}
{\( 115^{\circ} \)}{\( 120^{\circ} \)}{\( 90^{\circ} \)}{\( 105^{\circ} \)}{}{}}
\LANGpl{
\Question{
W przeszłości  w szkole nauki jazdy  uczono nas, że kierownicę należy trzymać w pozycji "za dziesięć  druga". Podaj miarę kąta między wskazówką godzinową a wskazówką minutową, która odpowiada tej godzinie.}
{\( 115^{\circ} \)}{\( 120^{\circ} \)}{\( 90^{\circ} \)}{\( 105^{\circ} \)}{}{}}
\LANGcs{
\Question{
V minulosti jsme se v autoškole učili, že volant se má držet v poloze „za deset dvě“. Vypočítejte velikost úhlu mezi hodinovou a minutovou ručičkou, které ukazují právě tento čas.}
{\( 115^{\circ} \)}{\( 120^{\circ} \)}{\( 90^{\circ} \)}{\( 105^{\circ} \)}{}{}}
\LANGsk{
\Question{
V minulosti sa v autoškole učilo, že volant treba držať v polohe „za desať dve“. Určte v stupňoch veľkosť uhla medzi hodinovou a minútovou ručičkou, ktorý odpovedá tomuto času.}
{\( 115^{\circ} \)}{\( 120^{\circ} \)}{\( 90^{\circ} \)}{\( 105^{\circ} \)}{}{}}
\LANGes{
\Question{
En el pasado en la autoescuela nos enseñaban a sujetar el volante en la posición "dos menos diez“. Cacula el ángulo entre la maniilla pequeña y la grande en ese momento.}
{\( 115^{\circ} \)}{\( 120^{\circ} \)}{\( 90^{\circ} \)}{\( 105^{\circ} \)}{}{}}
\End

\Data\ID{1003023306}
\Updated{2020-12-10 02:12:12}
\Level{B}
\SubArea{Angles, arcs and sectors}
\LANGen{
\Question{
The radian measure of the angle \( \theta \) is \( \frac{\pi}2 \). What is the sum of all the radian measures of the angles coterminal to \( \theta \) from the interval \( [-5\pi; 5\pi] \)? (Two angles are coterminal if they are drawn in the standard position and both have their terminal sides in the same location.)}
{\( \frac52\pi \)}{\( 0 \)}{\( 3\pi \)}{\( \pi \)}{}{}}
\LANGpl{
\Question{
Miara kąta \( \theta \) jest równa \( \frac{\pi}2 \). Jaka jest suma miar kątów równych kątowi \( \theta \) zawartych w przedziale  \( \langle-5\pi; 5\pi\rangle \) ?}
{\( \frac52\pi \)}{\( 0 \)}{\( 3\pi \)}{\( \pi \)}{}{}}
\LANGcs{
\Question{
Základní velikost orientovaného úhlu \( \theta \) je \( \frac{\pi}2 \). Součet všech jeho velikostí, které leží v intervalu \( \langle-5\pi; 5\pi\rangle \) je:}
{\( \frac52\pi \)}{\( 0 \)}{\( 3\pi \)}{\( \pi \)}{}{}}
\LANGsk{
\Question{
Základná veľkosť orientovaného uhla  \( \theta \) je \( \frac{\pi}2 \). Súčet všetkých jeho veľkostí, ktoré ležia  v intervale \( \langle-5\pi; 5\pi\rangle \)  je:}
{\( \frac52\pi \)}{\( 0 \)}{\( 3\pi \)}{\( \pi \)}{}{}}
\LANGes{
\Question{
La medida del ángulo \( \theta \) en radianes es \( \frac{\pi}2 \). ¿Qué es la suma de todas las medidas de los ángulos coterminales con \( \theta \) en radianes del intervalo \( [-5\pi; 5\pi] \)? (Dos ángulos son coterminales si dibujándolos en la posición estándar tienen los lados terminales en el mismo sitio.)}
{\( \frac52\pi \)}{\( 0 \)}{\( 3\pi \)}{\( \pi \)}{}{}}
\End

\Data\ID{1003023307}
\Updated{2020-12-10 03:12:26}
\Level{B}
\SubArea{Angles, arcs and sectors}
\LANGen{
\Question{
The radian measure of the angle \( \theta \) is \( \frac{\pi}3 \). How many values from the set 
\( M = \left\{\frac43\pi; \frac73\pi; \frac83\pi; \frac{13}3\pi; -\frac53\pi; \frac{61}3\pi; -\frac{61}3\pi \right\} \) are the measures of angles coterminal to \( \theta \)? (Two angles are coterminal if they are drawn in the standard position and both have their terminal sides in the same location.)}
{\( 4 \)}{\( 5 \)}{\( 6 \)}{\( 3 \)}{}{}}
\LANGpl{
\Question{
Miara kąta \( \theta \) jest równa \( \frac{\pi}3 \). Ile wartości ze zbioru 
\( M = \left\{\frac43\pi; \frac73\pi; \frac83\pi; \frac{13}3\pi; -\frac53\pi; \frac{61}3\pi; -\frac{61}3\pi \right\} \) jest miarą kąta skierowanego równą kątowi \( \theta \)?}
{\( 4 \)}{\( 5 \)}{\( 6 \)}{\( 3 \)}{}{}}
\LANGcs{
\Question{
Základní velikost orientovaného úhlu \( \theta \) je \( \frac{\pi}3 \). Kolik čísel z množiny M vyjadřuje velikost tohoto úhlu, jestliže \( M = \left\{\frac43\pi; \frac73\pi; \frac83\pi; \frac{13}3\pi; -\frac53\pi; \frac{61}3\pi; -\frac{61}3\pi \right\} \)?}
{\( 4 \)}{\( 5 \)}{\( 6 \)}{\( 3 \)}{}{}}
\LANGsk{
\Question{
Základná veľkosť orientovaného uhla \( \theta \) je \( \frac{\pi}3 \).  Koľko čísel z množiny M vyjadruje veľkosť tohto uhla, ak \( M = \left\{\frac43\pi; \frac73\pi; \frac83\pi; \frac{13}3\pi; -\frac53\pi; \frac{61}3\pi; -\frac{61}3\pi \right\} \) ?}
{\( 4 \)}{\( 5 \)}{\( 6 \)}{\( 3 \)}{}{}}
\LANGes{
\Question{
La medida del ángulo \( \theta \) en radianes es \( \frac{\pi}3 \). ¿Cuántos valores del conjunto
\( M = \left\{\frac43\pi; \frac73\pi; \frac83\pi; \frac{13}3\pi; -\frac53\pi; \frac{61}3\pi; -\frac{61}3\pi \right\} \) son medidas de ángulos coterminales a \( \theta \)? (Dos ángulos son coterminales si dibujándolos en la posición estándar tienen los lados terminales en el mismo sitio.)}
{\( 4 \)}{\( 5 \)}{\( 6 \)}{\( 3 \)}{}{}}
\End

\Data\ID{1003023308}
\Updated{2021-06-21 09:06:01}
\Level{B}
\SubArea{Angles, arcs and sectors}
\LANGen{
\Question{
What angle sweeps the minute hand of a clock in \( 10 \) minutes?}
{\( -60^{\circ} \)}{\( 36^{\circ} \)}{\( -36^{\circ} \)}{\( 60^{\circ} \)}{}{}}
\LANGpl{
\Question{
Jaki kąt łukowy tworzą wskazówki zegara w ciągu  \( 10 \) minut?}
{\( -60^{\circ} \)}{\( 36^{\circ} \)}{\( -36^{\circ} \)}{\( 60^{\circ} \)}{}{}}
\LANGcs{
\Question{
Jaká je základní velikost orientovaného úhlu, který opíše minutová ručička za \( 10 \) minut?}
{\( -60^{\circ} \)}{\( 36^{\circ} \)}{\( -36^{\circ} \)}{\( 60^{\circ} \)}{}{}}
\LANGsk{
\Question{
Aká je základná veľkosť orientovaného uhla, ktorý opíše minútová ručička za  \( 10 \) 

 



minút?}
{\( -60^{\circ} \)}{\( 36^{\circ} \)}{\( -36^{\circ} \)}{\( 60^{\circ} \)}{}{}}
\LANGes{
\Question{
¿Qué ángulo recorre la manecilla grande de un reloj a los \( 10 \) minutos?}
{\( -60^{\circ} \)}{\( 36^{\circ} \)}{\( -36^{\circ} \)}{\( 60^{\circ} \)}{}{}}
\End

\Data\ID{1003023309}
\Updated{2021-06-21 09:06:25}
\Level{B}
\SubArea{Angles, arcs and sectors}
\LANGen{
\Question{
What angle sweeps the hour hand of a clock in \( 3 \) hours?}
{\( -90^{\circ} \)}{\( 45^{\circ} \)}{\( -45^{\circ} \)}{\( 90^{\circ} \)}{}{}}
\LANGpl{
\Question{
Jaki kąt łukowy tworzą wskazówki zegara w ciągu \( 3 \) godzin?}
{\( -90^{\circ} \)}{\( 45^{\circ} \)}{\( -45^{\circ} \)}{\( 90^{\circ} \)}{}{}}
\LANGcs{
\Question{
Jaká je základní velikost orientovaného úhlu, který opíše malá hodinová ručička za \( 3 \) hodiny?}
{\( -90^{\circ} \)}{\( 45^{\circ} \)}{\( -45^{\circ} \)}{\( 90^{\circ} \)}{}{}}
\LANGsk{
\Question{
Aká je základná veľkosť orientovaného uhla, ktorý opíše malá hodinová ručička za  \( 3 \) hodiny?}
{\( -90^{\circ} \)}{\( 45^{\circ} \)}{\( -45^{\circ} \)}{\( 90^{\circ} \)}{}{}}
\LANGes{
\Question{
¿Qué ángulo recorre la manecilla pequeña de un reloj en \( 3 \) horas?}
{\( -90^{\circ} \)}{\( 45^{\circ} \)}{\( -45^{\circ} \)}{\( 90^{\circ} \)}{}{}}
\End

\Data\ID{1003055101}
\Updated{2020-12-15 09:12:01}
\Level{B}
\SubArea{Angles, arcs and sectors}
\LANGen{
\Question{
How many different values of \( \theta \) meet both conditions:
\[\theta\in\bigcup\limits_{k\in\mathbb{Z}}\left\{\frac{3\pi}4+2k\pi\right\}\text{, }\ \theta\in[-2\pi;4\pi ] \]}
{\( 3 \)}{\( 4 \)}{\( 5 \)}{\( 2 \)}{}{}}
\LANGpl{
\Question{
Ile różnych wartości  \( \theta \) spełnia oba warunki:
\[\theta\in\bigcup\limits_{k\in\mathbb{Z}}\left\{\frac{3\pi}4+2k\pi\right\}\text{, }\ \theta\in\langle-2\pi;4\pi \rangle\]}
{\( 3 \)}{\( 4 \)}{\( 5 \)}{\( 2 \)}{}{}}
\LANGcs{
\Question{
Kolik různých hodnot $\theta$ splňuje obě podmínky:
\[\theta\in\bigcup\limits_{k\in\mathbb{Z}}\left\{\frac{3\pi}4+2k\pi\right\}\text{, }\ \theta\in\langle-2\pi;4\pi \rangle\]}
{\( 3 \)}{\( 4 \)}{\( 5 \)}{\( 2 \)}{}{}}
\LANGsk{
\Question{
Koľko rôznych hodnôt $\theta$ spĺňa obe podmienky:
\[\theta\in\bigcup\limits_{k\in\mathbb{Z}}\left\{\frac{3\pi}4+2k\pi\right\}\text{, }\ \theta\in\langle-2\pi;4\pi \rangle \]}
{\( 3 \)}{\( 4 \)}{\( 5 \)}{\( 2 \)}{}{}}
\LANGes{
\Question{
¿Cuántas valores diferentes de \( \theta \) satisfacen ambas condiciones?
\[\theta\in\bigcup\limits_{k\in\mathbb{Z}}\left\{\frac{3\pi}4+2k\pi\right\}\text{, }\ \theta\in[-2\pi;4\pi ] \]}
{\( 3 \)}{\( 4 \)}{\( 5 \)}{\( 2 \)}{}{}}
\End

\Data\ID{1003055102}
\Updated{2021-06-21 09:06:36}
\Level{B}
\SubArea{Angles, arcs and sectors}
\LANGen{
\Question{
The measure of the angle \( \theta \) meets the following conditions:
\[\theta\in\bigcup\limits_{k\in\mathbb{Z}}\left\{\frac23\pi+k\frac{\pi}3\right\},\ \theta\in\left[-\frac{\pi}2;2\pi\right].\]
Choose the smallest value of $\theta$.}
{\( -\frac{\pi}3 \)}{\( -\frac{\pi}2 \)}{\( \frac23\pi \)}{\( \frac{\pi}3 \)}{}{}}
\LANGpl{
\Question{
Miary kąta \( \theta \) spełniają warunki:
\[\theta\in\bigcup\limits_{k\in\mathbb{Z}}\left\{\frac23\pi+k\frac{\pi}3\right\},\ \theta\in\left\langle-\frac{\pi}2;2\pi\right\rangle.\]
Wybierz najmniejszą wartość $\theta$.}
{\( -\frac{\pi}3 \)}{\( -\frac{\pi}2 \)}{\( \frac23\pi \)}{\( \frac{\pi}3 \)}{}{}}
\LANGcs{
\Question{
Velikost úhlu $\theta$ splňuje tyto podmínky:
\[\theta\in\bigcup\limits_{k\in\mathbb{Z}}\left\{\frac23\pi+k\frac{\pi}3\right\},\ \theta\in\left\langle-\frac{\pi}2;2\pi\right\rangle.\]
Vyberte nejmenší hodnotu $\theta$.}
{\( -\frac{\pi}3 \)}{\( -\frac{\pi}2 \)}{\( \frac23\pi \)}{\( \frac{\pi}3 \)}{}{}}
\LANGsk{
\Question{
Veľkosť uhla $\theta$ spĺňa tieto podmienky:
\[\theta\in\bigcup\limits_{k\in\mathbb{Z}}\left\{\frac23\pi+k\frac{\pi}3\right\},\ \theta\in\left\langle-\frac{\pi}2;2\pi\right\rangle.\]
Vyberte najmenšiu hodnotu $\theta$.}
{\( -\frac{\pi}3 \)}{\( -\frac{\pi}2 \)}{\( \frac23\pi \)}{\( \frac{\pi}3 \)}{}{}}
\LANGes{
\Question{
La medida del ángulo $\theta$ satisface estas condiciones:
\[\theta\in\bigcup\limits_{k\in\mathbb{Z}}\left\{\frac23\pi+k\frac{\pi}3\right\},\ \theta\in\left\langle-\frac{\pi}2;2\pi\right\rangle.\]
Elige el menor valor de $\theta$.}
{\( -\frac{\pi}3 \)}{\( -\frac{\pi}2 \)}{\( \frac23\pi \)}{\( \frac{\pi}3 \)}{}{}}
\End

\Data\ID{1003055103}
\Updated{2021-06-21 09:06:29}
\Level{B}
\SubArea{Angles, arcs and sectors}
\LANGen{
\Question{
Find out for which \( k\in\mathbb{Z} \) is true:
\[ -\frac{11}4\pi = \frac74\pi+k\frac{\pi}2 \]}
{\( -9 \)}{\( 9 \)}{\( -18 \)}{\( 18 \)}{}{}}
\LANGpl{
\Question{
Sprawdź, dla którego \( k\in\mathbb{Z} \) równanie jest prawdziwe:
\[ -\frac{11}4\pi = \frac74\pi+k\frac{\pi}2 \]}
{\( -9 \)}{\( 9 \)}{\( -18 \)}{\( 18 \)}{}{}}
\LANGcs{
\Question{
Zjistěte, pro které $k\in \mathbb{Z}$ platí:
\[ -\frac{11}4\pi = \frac74\pi+k\frac{\pi}2 \]}
{\( -9 \)}{\( 9 \)}{\( -18 \)}{\( 18 \)}{}{}}
\LANGsk{
\Question{
Zistite, pre ktoré $k\in \mathbb{Z}$ platí:
\[ -\frac{11}4\pi = \frac74\pi+k\frac{\pi}2 \]}
{\( -9 \)}{\( 9 \)}{\( -18 \)}{\( 18 \)}{}{}}
\LANGes{
\Question{
Averigua para cuáles \( k\in\mathbb{Z} \) es válida la igualdad :
\[ -\frac{11}4\pi = \frac74\pi+k\frac{\pi}2 \]}
{\( -9 \)}{\( 9 \)}{\( -18 \)}{\( 18 \)}{}{}}
\End

\Data\ID{1003055104}
\Updated{2021-01-04 06:01:19}
\Level{B}
\SubArea{Angles, arcs and sectors}
\LANGen{
\Question{
The measures of given angles belong to the set \( M \):
\[ M = \left\{ 120^{\circ} + k\,360^{\circ}\right\},\ k\in\{-1;0;1\}.\]
Determine their arithmetic mean.}
{\( 120^{\circ} \)}{\( 420^{\circ} \)}{\( -120^{\circ} \)}{\( -420^{\circ} \)}{}{}}
\LANGpl{
\Question{
Miary kątów należą do zbioru \( M \):
\[ M = \left\{ 120^{\circ} + k\,360^{\circ}\right\},\ k\in\{-1;0;1\}.\]
Podaj ich wartość arytmetyczną.}
{\( 120^{\circ} \)}{\( 420^{\circ} \)}{\( -120^{\circ} \)}{\( -420^{\circ} \)}{}{}}
\LANGcs{
\Question{
Hodnoty velikostí úhlů patří do množiny $M$:  
\[ M = \left\{ 120^{\circ} + k\,360^{\circ}\right\},\ k\in\{-1;0;1\}.\]
Určete jejich aritmetický průměr.}
{\( 120^{\circ} \)}{\( 420^{\circ} \)}{\( -120^{\circ} \)}{\( -420^{\circ} \)}{}{}}
\LANGsk{
\Question{
Hodnoty veľkostí uhlov patria do množiny $M$:  
\[ M = \left\{ 120^{\circ} + k\,360^{\circ}\right\},\ k\in\{-1;0;1\}.\]
Určte ich aritmetický priemer.}
{\( 120^{\circ} \)}{\( 420^{\circ} \)}{\( -120^{\circ} \)}{\( -420^{\circ} \)}{}{}}
\LANGes{
\Question{
Las medidas de los ángulos perteneces al conjunto $M$:  
\[ M = \left\{ 120^{\circ} + k\,360^{\circ}\right\},\ k\in\{-1;0;1\}.\]
Halla su media aritmética.}
{\( 120^{\circ} \)}{\( 420^{\circ} \)}{\( -120^{\circ} \)}{\( -420^{\circ} \)}{}{}}
\End

\Data\ID{1003055105}
\Updated{2021-01-04 06:01:28}
\Level{B}
\SubArea{Angles, arcs and sectors}
\LANGen{
\Question{
The measurements of angles \( \theta \) form  the set \( M \):
\[ M=\left\{\pi+k\frac23\pi\right\},\ k\in\{-2;-1;1;2\}.\]
Calculate their sum.}
{\( 4\pi \)}{\( 0 \)}{\( \frac23 \pi \)}{\( -\frac23 \pi \)}{}{}}
\LANGpl{
\Question{
Miary kątów \( \theta \) tworzą zbiór \( M \):
\[ M=\left\{\pi+k\frac23\pi\right\},\ k\in\{-2;-1;1;2\}.\]
Oblicz ich sumę.}
{\( 4\pi \)}{\( 0 \)}{\( \frac23 \pi \)}{\( -\frac23 \pi \)}{}{}}
\LANGcs{
\Question{
Velikosti úhlů \( \theta \) tvoří množinu $M$: 
\[ M=\left\{\pi+k\frac23\pi\right\},\ k\in\{-2;-1;1;2\}.\]
Určete jejich součet.}
{\( 4\pi \)}{\( 0 \)}{\( \frac23 \pi \)}{\( -\frac23 \pi \)}{}{}}
\LANGsk{
\Question{
Veľkosti uhlov \( \theta \)  tvoria množinu $M$: 
\[ M=\left\{\pi+k\frac23\pi\right\},\ k\in\{-2;-1;1;2\}.\]
Určte ich súčet.}
{\( 4\pi \)}{\( 0 \)}{\( \frac23 \pi \)}{\( -\frac23 \pi \)}{}{}}
\LANGes{
\Question{
Las medidas de los ángulos \( \theta \) construyen el conjunto $M$: 
\[ M=\left\{\pi+k\frac23\pi\right\},\ k\in\{-2;-1;1;2\}.\]
Halla su suma.}
{\( 4\pi \)}{\( 0 \)}{\( \frac23 \pi \)}{\( -\frac23 \pi \)}{}{}}
\End

\Data\ID{1003055106}
\Updated{2021-01-04 06:01:21}
\Level{B}
\SubArea{Angles, arcs and sectors}
\LANGen{
\Question{
Given two angles \( \alpha=  \frac{66}{15}\pi \) and \( \beta=\frac{*}{5}\pi \), which of the following numbers is to be substituted for \( * \) so that both angles  take the same position on the unit circle?}
{\( 2 \)}{\( 1 \)}{\( 0 \)}{\( 3 \)}{}{}}
\LANGpl{
\Question{
Dane są dwa kąty \( \alpha=  \frac{66}{15}\pi \) i \( \beta=\frac{*}{5}\pi \),  którą liczbę należy wstawić w miejsce  \( * \), aby dwa kąty miały taką samą pozycję na kole jednostkowym?}
{\( 2 \)}{\( 1 \)}{\( 0 \)}{\( 3 \)}{}{}}
\LANGcs{
\Question{
Jsou dány dva úhly \( \alpha=  \frac{66}{15}\pi \) a \( \beta=\frac{*}{5}\pi \). Které z následujících čísel je třeba dosadit na místo \( * \), aby se dané úhly zobrazily na jednotkové kružnici stejně?}
{\( 2 \)}{\( 1 \)}{\( 0 \)}{\( 3 \)}{}{}}
\LANGsk{
\Question{
Dané sú dva uhly  \( \alpha=  \frac{66}{15}\pi \) and \( \beta=\frac{*}{5}\pi \). Ktorým z nasledujúcich čísel treba nahradiť \( * \), aby sa dané uhly zobrazili na jednotkovej kružnici rovnako?}
{\( 2 \)}{\( 1 \)}{\( 0 \)}{\( 3 \)}{}{}}
\LANGes{
\Question{
Sean dos ángulos \( \alpha=  \frac{66}{15}\pi \) y \( \beta=\frac{*}{5}\pi \). ¿Cuál de los números tiene que estar en vez de \( * \) para que los ángulos aparezcan en el mismo lugar en la circunferencia unitaria?}
{\( 2 \)}{\( 1 \)}{\( 0 \)}{\( 3 \)}{}{}}
\End

\Data\ID{1003055109}
\Updated{2020-12-13 09:12:30}
\Level{B}
\SubArea{Angles, arcs and sectors}
\LANGen{
\Question{
What is the maximum error that can be obtained by adding four angles, if the size of each of them is rounded to the nearest degree before the addition?}
{\( 2^{\circ} \)}{\( 0^{\circ} \)}{\( 3^{\circ} \)}{\( 4^{\circ} \)}{}{}}
\LANGpl{
\Question{
Jaki jest maksymalny błąd, który można zrobić, dodając cztery kąty, jeśli wielkość każdego z nich jest zaokrąglana do pełnego stopnia przed dodaniem?}
{\( 2^{\circ} \)}{\( 0^{\circ} \)}{\( 3^{\circ} \)}{\( 4^{\circ} \)}{}{}}
\LANGcs{
\Question{
Jakou největší chybu můžeme způsobit při sčítání čtyř úhlů, jestliže velikost každého z nich zaokrouhlíme na celé stupně ještě před sčítáním?}
{\( 2^{\circ} \)}{\( 0^{\circ} \)}{\( 3^{\circ} \)}{\( 4^{\circ} \)}{}{}}
\LANGsk{
\Question{
Akú najväčšiu nepresnosť môžeme dosiahnuť pri sčítaní štyroch uhlov, ak veľkosť každého z nich zaokrúhlime na celé stupne ešte pred sčítaním?}
{\( 2^{\circ} \)}{\( 0^{\circ} \)}{\( 3^{\circ} \)}{\( 4^{\circ} \)}{}{}}
\LANGes{
\Question{
¿Qué es el error máximo que podemos obtener sumando cuatro ángulos, si cada uno de ellos está aproximado al grado más cercano antes de sumar?}
{\( 2^{\circ} \)}{\( 0^{\circ} \)}{\( 3^{\circ} \)}{\( 4^{\circ} \)}{}{}}
\End

\Data\ID{1003055110}
\Updated{2021-01-06 07:01:24}
\Level{B}
\SubArea{Angles, arcs and sectors}
\LANGen{
\Question{
How many times in \( 12 \) hours (from \( 0\!:\!00 \) to \( 11\!:\!59\!:\!59 \)) will the minute hand and the hour hand make an angle of \( 1 \) degree?}
{\( 22 \)}{\( 11 \)}{\( 12 \)}{\( 24 \)}{}{}}
\LANGpl{
\Question{
Ile razy w ciągu \( 12 \) godzin (od \( 0\!:\!00 \) do \( 11\!:\!59\!:\!59 \)) wskazówka minutowa i wskazówka godzinowa utworzą kąt \( 1 \) stopnia?}
{\( 22 \)}{\( 11 \)}{\( 12 \)}{\( 24 \)}{}{}}
\LANGcs{
\Question{
Kolikrát v průběhu 12 hodin (počínaje \( 0\!:\!00 \)  a konče \( 11\!:\!59\!:\!59 \)) budou ručičky hodin svírat úhel \( 1 \) stupeň?}
{\( 22 \)}{\( 11 \)}{\( 12 \)}{\( 24 \)}{}{}}
\LANGsk{
\Question{
Koľkokrát v priebehu 12 hodín (počnúc \( 0\!:\!00 \)  a končiac  \( 11\!:\!59\!:\!59 \)) budú ručičky hodín   zvierať uhol \( 1 \) stupeň?}
{\( 22 \)}{\( 11 \)}{\( 12 \)}{\( 24 \)}{}{}}
\LANGes{
\Question{
¿Cuántas veces las manillas de horas forman el ángulo de \( 1 \) grado durante 12 horas (desde \( 0\!:\!00 \)  hasta \( 11\!:\!59\!:\!59 \))?}
{\( 22 \)}{\( 11 \)}{\( 12 \)}{\( 24 \)}{}{}}
\End

\Data\ID{1003055201}
\Updated{2021-01-06 07:01:05}
\Level{B}
\SubArea{Angles, arcs and sectors}
\LANGen{
\Question{
In \( 12 \) hours (from \( 0\!:\!00 \) am to \( 11\!:\!59\!:\!59 \) am), how many times do the minute and the hour hand make an angle of \( 0 \) rad ?}
{\( 11 \)}{\( 22 \)}{\( 12 \)}{\( 24 \)}{}{}}
\LANGpl{
\Question{
Ile razy w ciągu  \( 12 \) godzin (od \( 0\!:\!00 \)  do \( 11\!:\!59\!:\!59 \)), wskazówka minutowa i godzinowa utworzą  kąt \( 0 \) rad ?}
{\( 11 \)}{\( 22 \)}{\( 12 \)}{\( 24 \)}{}{}}
\LANGcs{
\Question{
Kolikrát v průběhu \( 12 \) hodin  (počínaje \( 0\!:\!00 \) a konče \( 11\!:\!59\!:\!59 \)), budou ručičky hodin svírat úhel \( 0 \) rad ?}
{\( 11 \)}{\( 22 \)}{\( 12 \)}{\( 24 \)}{}{}}
\LANGsk{
\Question{
Koľkokrát v priebehu \( 12 \) hodín  (počnúc \( 0\!:\!00 \) a končiac \( 11\!:\!59\!:\!59 \)), budú ručičky hodín zvierať uhol \( 0 \) rad ?}
{\( 11 \)}{\( 22 \)}{\( 12 \)}{\( 24 \)}{}{}}
\LANGes{
\Question{
¿Cuántas veces las manillas de horas forman el ángulo de  \( 0 \) radianes durante \( 12 \) horas  (desde \( 0\!:\!00 \) hasta \( 11\!:\!59\!:\!59 \))?}
{\( 11 \)}{\( 22 \)}{\( 12 \)}{\( 24 \)}{}{}}
\End

\Data\ID{1003055202}
\Updated{2021-01-06 07:01:36}
\Level{B}
\SubArea{Angles, arcs and sectors}
\LANGen{
\Question{
Choose a statement that is not true  for the angle \( \theta \) that the minute hand as an initial arm makes with the hour hand as a terminal arm in the clockwise direction.}
{\( 4\!:\!30 \), \( \theta = -300^{\circ}30' \)}{\( 10\!:\!45 \), \( \theta = -52^{\circ}30' \)}{\( 7\!:\!45\), \( \theta = -322^{\circ}30' \)}{\( 3\!:\!15 \), \( \theta = -7^{\circ}30' \)}{}{}}
\LANGpl{
\Question{
Wskaż zależność, która nie ma zastosowania do kąta  \( \theta \), którego wskazówka minutowa to ramię początkowe, a wskazówka godzinowa to ramię końcowe zgodnie ze kierunkiem ruchu zegara.}
{\( 4\!:\!30 \), \( \theta = -300^{\circ}30' \)}{\( 10\!:\!45 \), \( \theta = -52^{\circ}30' \)}{\( 7\!:\!45\), \( \theta = -322^{\circ}30' \)}{\( 3\!:\!15 \), \( \theta = -7^{\circ}30' \)}{}{}}
\LANGcs{
\Question{
Vyberte nepravdivé tvrzení, které neplatí pro úhel \( \theta \), který svírá velká hodinová ručička jako počáteční rameno a malá hodinová ručička jako koncové rameno úhlu ve směru hodinových ručiček.}
{\( 4\!:\!30 \), \( \theta = -300^{\circ}30' \)}{\( 10\!:\!45 \), \( \theta = -52^{\circ}30' \)}{\( 7\!:\!45\), \( \theta = -322^{\circ}30' \)}{\( 3\!:\!15 \), \( \theta = -7^{\circ}30' \)}{}{}}
\LANGsk{
\Question{
Vyberte nepravdivé tvrdenie, ktoré neplatí pre uhol \( \theta \), ktorý zviera veľká hodinová ručička ako začiatočné rameno a malá hodinová ručička ako koncové rameno uhla v zmysle chodu hodín.}
{\( 4\!:\!30 \), \( \theta = -300^{\circ}30' \)}{\( 10\!:\!45 \), \( \theta = -52^{\circ}30' \)}{\( 7\!:\!45\), \( \theta = -322^{\circ}30' \)}{\( 3\!:\!15 \), \( \theta = -7^{\circ}30' \)}{}{}}
\LANGes{
\Question{
Elije la declaración falsa para el ángulo \( \theta \) que forma la manilla grande como lado inicial con la manilla pequeña como lado terminal en dirección horaria.}
{\( 4\!:\!30 \), \( \theta = -300^{\circ}30' \)}{\( 10\!:\!45 \), \( \theta = -52^{\circ}30' \)}{\( 7\!:\!45\), \( \theta = -322^{\circ}30' \)}{\( 3\!:\!15 \), \( \theta = -7^{\circ}30' \)}{}{}}
\End

\Data\ID{1003055203}
\Updated{2021-01-13 11:01:53}
\Level{B}
\SubArea{Angles, arcs and sectors}
\LANGen{
\Question{
Let the minute hand determine the initial arm and let the hour hand determine the terminal arm of an angle in the clockwise direction. What is the radian measure of the angle between the two hands at \( 11\!:\!30 \)?}
{\( -\frac{11}{12}\pi \)}{\( -\frac56\pi \)}{\( -\frac76\pi \)}{\( -\frac{13}{12}\pi \)}{}{}}
\LANGpl{
\Question{
Wskazówka minutowa zegara  to  ramię początkowe kąta, wskazówka godzinowa to ramię końcowe kąta,  kierunek zgodny z ruchem wskazówek zegara. Jaka jest miara kąta w radianach pomiędzy dwoma wskazówkami o godzinie \( 11\!:\!30 \)?}
{\( -\frac{11}{12}\pi \)}{\( -\frac56\pi \)}{\( -\frac76\pi \)}{\( -\frac{13}{12}\pi \)}{}{}}
\LANGcs{
\Question{
Jestliže velká hodinová ručička představuje počáteční rameno a malá hodinová ručička koncové rameno úhlu ve směru pohybu hodinových ručiček, jaký úhel v radiánech svírají ručičky v \( 11\!:\!30 \)?}
{\( -\frac{11}{12}\pi \)}{\( -\frac56\pi \)}{\( -\frac76\pi \)}{\( -\frac{13}{12}\pi \)}{}{}}
\LANGsk{
\Question{
Ak veľká hodinová ručička určuje začiatočné rameno a malá hodinová ručička koncové rameno uhla v zmysle chodu hodín, aký uhol v radiánoch zvierajú ručičky o \( 11\!:\!30 \)?}
{\( -\frac{11}{12}\pi \)}{\( -\frac56\pi \)}{\( -\frac76\pi \)}{\( -\frac{13}{12}\pi \)}{}{}}
\LANGes{
\Question{
¿Si la manilla grande es el lado inicial y la manilla pequeña es el lado terminal de un ángulo en la dirección horaria, qué medida en radianes tiene el ángulo en \( 11\!:\!30 \)?}
{\( -\frac{11}{12}\pi \)}{\( -\frac56\pi \)}{\( -\frac76\pi \)}{\( -\frac{13}{12}\pi \)}{}{}}
\End

\Data\ID{1003055204}
\Updated{2021-01-13 11:01:17}
\Level{B}
\SubArea{Angles, arcs and sectors}
\LANGen{
\Question{
If the minute hand determines the initial arm and if the hour hand determines the terminal arm of an angle in the clockwise direction, what is the radian measure of the angle between the two hands at \( 5\!:\!00 \)?}
{\( -\frac56\pi \)}{\( -\frac76\pi \)}{\( -\frac34\pi \)}{\( -\frac23\pi \)}{}{}}
\LANGpl{
\Question{
Jeśli wskazówka minutowa określa ramię początkowe a wskazówka godzinowa określa ramię końcowe kąta w kierunku zgodnym z ruchem wskazówek zegara, to jaka jest  miara kąta w radianach między między dwoma wskazówkami zegara o godzinie  \( 5\!:\!00 \)?}
{\( -\frac56\pi \)}{\( -\frac76\pi \)}{\( -\frac34\pi \)}{\( -\frac23\pi \)}{}{}}
\LANGcs{
\Question{
Jestliže velká hodinová ručička představuje počáteční rameno a malá hodinová ručička koncové rameno úhlu ve směru pohybu hodinových ručiček, jaký úhel v radiánech svírají ručičky v \( 5\!:\!00 \)?}
{\( -\frac56\pi \)}{\( -\frac76\pi \)}{\( -\frac34\pi \)}{\( -\frac23\pi \)}{}{}}
\LANGsk{
\Question{
Ak veľká hodinová ručička určuje začiatočné rameno a malá hodinová ručička koncové rameno uhla v zmysle chodu hodín, aký uhol v radiánoch zvierajú ručičky o \( 5\!:\!00 \)?}
{\( -\frac56\pi \)}{\( -\frac76\pi \)}{\( -\frac34\pi \)}{\( -\frac23\pi \)}{}{}}
\LANGes{
\Question{
¿Si la manilla grande es el lado inicial y la manilla pequeña es el lado terminal de un ángulo en la dirección horaria, qué medida en radianes tiene el ángulo en  \( 5\!:\!00 \)?}
{\( -\frac56\pi \)}{\( -\frac76\pi \)}{\( -\frac34\pi \)}{\( -\frac23\pi \)}{}{}}
\End

\Data\ID{1003055207}
\Updated{2021-01-13 12:01:56}
\Level{B}
\SubArea{Angles, arcs and sectors}
\LANGen{
\Question{
The measure of the angle \( \theta \) is \( \frac{\pi}4 \). How many coterminal angles to \( \theta \) have the measure from  the interval \( [ -4\pi;6\pi ] \)? (Two angles are coterminal if they are drawn in the standard position and both have their terminal sides in the same location.)}
{\( 5 \)}{\( 6 \)}{\( 7 \)}{\( 8 \)}{}{}}
\LANGpl{
\Question{
Miara kąta skierowanego   \( \theta \)  to \( \frac{\pi}4 \). Ile miar kątów \( \theta \)  jest w przedziale \( \langle -4\pi;6\pi \rangle \)?}
{\( 5 \)}{\( 6 \)}{\( 7 \)}{\( 8 \)}{}{}}
\LANGcs{
\Question{
Základní velikost orientovaného úhlu \( \theta \) je \( \frac{\pi}4 \). Kolik je možných velikostí úhlu \( \theta \) v intervalu \( \langle -4\pi;6\pi \rangle \)?}
{\( 5 \)}{\( 6 \)}{\( 7 \)}{\( 8 \)}{}{}}
\LANGsk{
\Question{
Základná veľkosť orientovaného uhla \( \theta \) je \( \frac{\pi}4 \). Koľko je všetkých veľkostí uhla \( \theta \) v intervale \( \langle -4\pi;6\pi \rangle \)?}
{\( 5 \)}{\( 6 \)}{\( 7 \)}{\( 8 \)}{}{}}
\LANGes{
\Question{
La medida canónica de l ángulo \( \theta \) es \( \frac{\pi}4 \). ¿Cuántas medidas posibles del ángulo \( \theta \) hay en el intervalo \( \langle -4\pi;6\pi \rangle \)?}
{\( 5 \)}{\( 6 \)}{\( 7 \)}{\( 8 \)}{}{}}
\End

\Data\ID{1003055208}
\Updated{2021-01-04 07:01:51}
\Level{B}
\SubArea{Angles, arcs and sectors}
\LANGen{
\Question{
Let \( k\in\{-5; - 4; -3; -2\} \). Which number \( k \) satisfies the equation:    \( -\frac{32}7 \pi=k\pi+\frac k7\pi \)?}
{\( -4 \)}{\( -5 \)}{\( -2 \)}{\( -3 \)}{}{}}
\LANGpl{
\Question{
Dano \( k\in\{-5; - 4; -3; -2\} \). Która liczba \( k \) spełnia równanie:    \( -\frac{32}7 \pi=k\pi+\frac k7\pi \)?}
{\( -4 \)}{\( -5 \)}{\( -2 \)}{\( -3 \)}{}{}}
\LANGcs{
\Question{
Které z čísel \( k\in\{-5; - 4; -3; -2\} \) vyhovuje rovnici: \( -\frac{32}7 \pi=k\pi+\frac k7\pi \)?}
{\( -4 \)}{\( -5 \)}{\( -2 \)}{\( -3 \)}{}{}}
\LANGsk{
\Question{
Ktoré z čísel \( k\in\{-5; - 4; -3; -2\} \) vyhovuje rovnici:   \( -\frac{32}7 \pi=k\pi+\frac k7\pi \)?}
{\( -4 \)}{\( -5 \)}{\( -2 \)}{\( -3 \)}{}{}}
\LANGes{
\Question{
¿Cuál de los números \( k\in\{-5; - 4; -3; -2\} \) satisface la ecuación: \( -\frac{32}7 \pi=k\pi+\frac k7\pi \)?}
{\( -4 \)}{\( -5 \)}{\( -2 \)}{\( -3 \)}{}{}}
\End

\Data\ID{1003055209}
\Updated{2021-01-04 07:01:20}
\Level{B}
\SubArea{Angles, arcs and sectors}
\LANGen{
\Question{
Let \( k\in\{5; 6; 7; 8\} \). Which number \( k \) satisfies the equation: \( \frac{91}{12}\pi=\pi k+\frac{\pi k}{12} \)?}
{\( 7 \)}{\( 8 \)}{\( 6 \)}{\( 5 \)}{}{}}
\LANGpl{
\Question{
Dane jest  \( k\in\{5; 6; 7; 8\} \). Jaka liczba \( k \) spełnia równanie: \( \frac{91}{12}\pi=\pi k+\frac{\pi k}{12} \)?}
{\( 7 \)}{\( 8 \)}{\( 6 \)}{\( 5 \)}{}{}}
\LANGcs{
\Question{
Které z čísel \( k\in\{5; 6; 7; 8\} \) vyhovuje rovnici \( \frac{91}{12}\pi=\pi k+\frac{\pi k}{12} \)?}
{\( 7 \)}{\( 8 \)}{\( 6 \)}{\( 5 \)}{}{}}
\LANGsk{
\Question{
Ktoré z čísel \( k\in\{5; 6; 7; 8\} \)  vyhovuje rovnici  \( \frac{91}{12}\pi=\pi k+\frac{\pi k}{12} \)?}
{\( 7 \)}{\( 8 \)}{\( 6 \)}{\( 5 \)}{}{}}
\LANGes{
\Question{
¿Cuál de los números \( k\in\{5; 6; 7; 8\} \) satisface la ecuación \( \frac{91}{12}\pi=\pi k+\frac{\pi k}{12} \)?}
{\( 7 \)}{\( 8 \)}{\( 6 \)}{\( 5 \)}{}{}}
\End

\Data\ID{1103023101}
\Updated{2021-01-13 12:01:35}
\Level{B}
\SubArea{Angles, arcs and sectors}
\IMG{1103023101.pdf}
\LANGen{
\Question{
Let \( PA \) be the initial side of an angle of \( 16.8 \) radians (see the picture). Which of the points \( K \), \( L \), \( M \), \( N \) lies on its terminal side?}
{\( M \)}{\( K \)}{\( L \)}{\( N \)}{}{}}
\LANGpl{
\Question{
Niech \( PA \) będzie ramieniem początkowym kąta  \( 16{,}8 \) radianów (spójrz na rysunek). Które z punktów \( K \), \( L \), \( M \), \( N \) leżą na jego ramieniu końcowym?}
{\( M \)}{\( K \)}{\( L \)}{\( N \)}{}{}}
\LANGcs{
\Question{
Nechť \( PA \) je počáteční rameno úhlu, jehož velikost v obloukové míře je \( 16{,}8 \). Který z bodů \( K \), \( L \), \( M \), \( N \)  leží na jeho koncovém rameni?}
{\( M \)}{\( K \)}{\( L \)}{\( N \)}{}{}}
\LANGsk{
\Question{
Nech  \( PA \) je počiatočné rameno uhla, ktorého veľkosť v  oblúkovej miere je  \( 16{,}8 \) . Ktorý z bodov \( K \), \( L \), \( M \), \( N \)  leží na jeho koncovom ramene?}
{\( M \)}{\( K \)}{\( L \)}{\( N \)}{}{}}
\LANGes{
\Question{
Sea \( PA \) el lado inicial del ángulo que mide \( 16{,}8 \) radianes. ¿Cuál de los puntos \( K \), \( L \), \( M \), \( N \)  pertenece a su lado terminal?}
{\( M \)}{\( K \)}{\( L \)}{\( N \)}{}{}}
\End

\Data\ID{1103023102}
\Updated{2021-01-13 12:01:40}
\Level{B}
\SubArea{Angles, arcs and sectors}
\IMG{1103023102.pdf}
\LANGen{
\Question{
Let \( PA \) be the initial side of an angle of \( -26.42 \) radians (see the picture). Which of the points \( K \), \( L \), \( M \), \( N \) lies on its terminal side?}
{\( N \)}{\( K \)}{\( L \)}{\( M \)}{}{}}
\LANGpl{
\Question{
Niech \( PA \) będzie ramieniem początkowym kąta \( -26{,}42 \) radianów (spójrz na rysunek). Które z punktów \( K \), \( L \), \( M \), \( N \) leżą na jego ramieniu końcowym?}
{\( N \)}{\( K \)}{\( L \)}{\( M \)}{}{}}
\LANGcs{
\Question{
Nechť \( PA \) je počáteční rameno úhlu, jehož velikost je \( -26{,}42 \)  rad. Který z bodů  \( K \), \( L \), \( M \), \( N \)  leží na jeho koncovém rameni?}
{\( N \)}{\( K \)}{\( L \)}{\( M \)}{}{}}
\LANGsk{
\Question{
Nech \( PA \) je počiatočné rameno uhla, ktorého veľkosť je \( -26{,}42 \)  rad. Ktorý z bodov  \( K \), \( L \), \( M \), \( N \)  leží na jeho koncovom ramene?}
{\( N \)}{\( K \)}{\( L \)}{\( M \)}{}{}}
\LANGes{
\Question{
Sea \( PA \) lado inicial del ángulo que mide \( -26{,}42 \)  rad. ¿Cuál de los puntos  \( K \), \( L \), \( M \), \( N \)  pertenece a su lado terminal?}
{\( N \)}{\( K \)}{\( L \)}{\( M \)}{}{}}
\End

\Data\ID{1103023103}
\Updated{2021-01-13 12:01:08}
\Level{B}
\SubArea{Angles, arcs and sectors}
\IMG{1103023103.pdf}
\LANGen{
\Question{
Let \( PA \) be the initial side of an angle of \( -\frac{17}3\pi \) radians (see the picture). Which of the points \( K \), \( L \), \( M \), \( N \) lies on its terminal side?}
{\( K \)}{\( L \)}{\( M \)}{\( N \)}{}{}}
\LANGpl{
\Question{
Niech \( PA \) będzie ramieniem początkowym kąta \( -\frac{17}3\pi \) radianów (spójrz na rysunek). Które z punktów \( K \), \( L \), \( M \), \( N \) leżą na jego ramieniu końcowym?}
{\( K \)}{\( L \)}{\( M \)}{\( N \)}{}{}}
\LANGcs{
\Question{
Nechť \( PA \) je počáteční rameno úhlu, jehož velikost v obloukové míře je \( -\frac{17}3\pi \). Který z bodů  \( K \), \( L \), \( M \), \( N \) leží na jeho koncovém rameni?}
{\( K \)}{\( L \)}{\( M \)}{\( N \)}{}{}}
\LANGsk{
\Question{
Nech \( PA \) je počiatočné rameno uhla, ktorého veľkosť v oblúkovej  miere je\( -\frac{17}3\pi \). Ktorý z bodov  \( K \), \( L \), \( M \), \( N \) leží na jeho koncovom ramene?}
{\( K \)}{\( L \)}{\( M \)}{\( N \)}{}{}}
\LANGes{
\Question{
Sea \( PA \) lado inicial del ángulo que mide \( -\frac{17}3\pi \). ¿Cuál de los puntos \( K \), \( L \), \( M \), \( N \) pertenece a su lado terminal?}
{\( K \)}{\( L \)}{\( M \)}{\( N \)}{}{}}
\End

\Data\ID{1103055107}
\Updated{2021-01-13 01:01:24}
\Level{B}
\SubArea{Angles, arcs and sectors}
\IMG{1103055107_8.pdf}
\LANGen{
\Question{
The figure shows a directional rosette that can be used to determine a marching angle. (The initial arm always faces north and the terminal arm determines the direction of the march, so the measure of the angle increases from north to east). Give the degree measure of the marching angle if the march is directed southwest.}
{\( 225^{\circ} \)}{\( 135^{\circ} \)}{\( -225^{\circ} \)}{\( -135^{\circ} \)}{}{}}
\LANGpl{
\Question{
Rysunek pokazuje kierunki świata, które mogą zostać użyte do określenia kąta. (Ramię początkowe wskazuje północ, ramię końcowe określa kierunek ruchu, zatem miara kąta wzrasta z północy na wschód.) 
Podaj miarę kąta jeśli ruch odbywa się na południowy zachód.}
{\( 225^{\circ} \)}{\( 135^{\circ} \)}{\( -225^{\circ} \)}{\( -135^{\circ} \)}{}{}}
\LANGcs{
\Question{
Na obrázku je kompas, pomocí kterého určujeme směr pochodu. (Počáteční rameno vždy směřuje na sever a koncové určuje směr pochodu, hodnoty tedy rostou od severu směrem k východu). 
Jak velký je pochodový úhel, jestliže je směr pochodu jihozápadní?}
{\( 225^{\circ} \)}{\( 135^{\circ} \)}{\( -225^{\circ} \)}{\( -135^{\circ} \)}{}{}}
\LANGsk{
\Question{
Na obrázku je smerová ružica, pomocou ktorej určujeme uhol pochodu. (Počiatočné rameno vždy smeruje na sever a koncové určuje smer pochodu, teda hodnoty pribúdajú od severu smerom k východu). 
Aký veľký je pochodový uhol, ak smer pochodu je juhozápadný?}
{\( 225^{\circ} \)}{\( 135^{\circ} \)}{\( -225^{\circ} \)}{\( -135^{\circ} \)}{}{}}
\LANGes{
\Question{
En el dibujo hay una brújula que se usa para averiguar la dirección del camino. (El lado inicial siempre señala al norte y el lado terminal señala la dirección del camino, es decir, las medidas crecen del norte hacia el este). ¿Qué medida tiene el ángulo de camino si la dirección del camino es suroeste?}
{\( 225^{\circ} \)}{\( 135^{\circ} \)}{\( -225^{\circ} \)}{\( -135^{\circ} \)}{}{}}
\End

\Data\ID{1103055108}
\Updated{2021-01-13 01:01:11}
\Level{B}
\SubArea{Angles, arcs and sectors}
\IMG{1103055107_8_0.pdf}
\LANGen{
\Question{
The figure shows a directional rosette that can be used to determine a marching angle. (The initial arm always faces north and the terminal arm determines the direction of the march, so the measure of the angle increases from north to east.) Give the radian measure of the marching angle if the march is directed  southeast.}
{\( \frac34 \pi \)}{\( \frac54 \pi \)}{\( -\frac34 \pi \)}{\( -\frac54 \pi \)}{}{}}
\LANGpl{
\Question{
Rysunek przedstawia kierunki geograficzne, które mogą być użyte do określenia ruchu kątowego. (Początkowe ramię zawsze skierowane jest na północ, a ramię końcowe określa kierunek ruchu, więc miara kąta wzrasta z północy na wschód.) Podaj miarę ruchu kąta w radianach, jeśli ruch jest skierowany na południowy wschód.}
{\( \frac34 \pi \)}{\( \frac54 \pi \)}{\( -\frac34 \pi \)}{\( -\frac54 \pi \)}{}{}}
\LANGcs{
\Question{
Na obrázku je kompas, pomocí kterého určujeme úhel pochodu. (Počáteční rameno vždy směřuje na sever a koncové určuje směr pochodu, hodnoty tedy rostou od severu směrem k východu.)
Jak velký je pochodový úhel, jestliže je směr pochodu jihovýchodní?}
{\( \frac34 \pi \)}{\( \frac54 \pi \)}{\( -\frac34 \pi \)}{\( -\frac54 \pi \)}{}{}}
\LANGsk{
\Question{
Na obrázku je smerová ružica, pomocou ktorej určujeme uhol pochodu. (Počiatočné rameno vždy smeruje na sever a koncové určuje smer pochodu,  teda hodnoty pribúdajú od severu smerom k východu.) 
Aký veľký je pochodový uhol, ak smer pochodu je juhovýchodný?}
{\( \frac34 \pi \)}{\( \frac54 \pi \)}{\( -\frac34 \pi \)}{\( -\frac54 \pi \)}{}{}}
\LANGes{
\Question{
En el dibujo hay una brújula que se usa para averiguar la dirección del camino. (El lado inicial siempre señala al norte y el lado terminal señala la dirección del camino, es decir, las medidas crecen del norte hacia el este). ¿Qué medida tiene el ángulo de camino si la dirección del camino es sureste?}
{\( \frac34 \pi \)}{\( \frac54 \pi \)}{\( -\frac34 \pi \)}{\( -\frac54 \pi \)}{}{}}
\End

\Data\ID{1103055205}
\Updated{2021-01-13 01:01:45}
\Level{B}
\SubArea{Angles, arcs and sectors}
\IMG{1103055205.pdf}
\LANGen{
\Question{
Given the square \( ABCD \), find the measures of all coterminal angles to the angle \( DCB \). (Two angles are coterminal if they are drawn in the standard position and both have their terminal sides in the same location.)}
{\( \frac{\pi}2+2k\pi \), \( k\in\mathbb{Z} \)}{\( \frac{\pi}2+k\pi \), \( k\in\mathbb{Z} \)}{\( -\frac{\pi}2+2k\pi \), \( k\in\mathbb{Z} \)}{\( -\frac{\pi}2+k\pi \), \( k\in\mathbb{Z} \)}{}{}}
\LANGpl{
\Question{
Dany jest kwadrat  \( ABCD \), wszystkie kąty skierowane na  \( DCB \) można zapisać jako:}
{\( \frac{\pi}2+2k\pi \), \( k\in\mathbb{Z} \)}{\( \frac{\pi}2+k\pi \), \( k\in\mathbb{Z} \)}{\( -\frac{\pi}2+2k\pi \), \( k\in\mathbb{Z} \)}{\( -\frac{\pi}2+k\pi \), \( k\in\mathbb{Z} \)}{}{}}
\LANGcs{
\Question{
Je dán čtverec \( ABCD \). Všechny velikosti orientovaného úhlu  \( DCB \) je možné zapsat ve tvaru:}
{\( \frac{\pi}2+2k\pi \), \( k\in\mathbb{Z} \)}{\( \frac{\pi}2+k\pi \), \( k\in\mathbb{Z} \)}{\( -\frac{\pi}2+2k\pi \), \( k\in\mathbb{Z} \)}{\( -\frac{\pi}2+k\pi \), \( k\in\mathbb{Z} \)}{}{}}
\LANGsk{
\Question{
Daný je štvorec \( ABCD \). Všetky veľkosti orientovaného uhla  \( DCB \) možno napísať v tvare:}
{\( \frac{\pi}2+2k\pi \), \( k\in\mathbb{Z} \)}{\( \frac{\pi}2+k\pi \), \( k\in\mathbb{Z} \)}{\( -\frac{\pi}2+2k\pi \), \( k\in\mathbb{Z} \)}{\( -\frac{\pi}2+k\pi \), \( k\in\mathbb{Z} \)}{}{}}
\LANGes{
\Question{
Sea un cuadrado \( ABCD \). Todas las medidas del ángulo orientado \( DCB \) se pueden escribir de forma:}
{\( \frac{\pi}2+2k\pi \), \( k\in\mathbb{Z} \)}{\( \frac{\pi}2+k\pi \), \( k\in\mathbb{Z} \)}{\( -\frac{\pi}2+2k\pi \), \( k\in\mathbb{Z} \)}{\( -\frac{\pi}2+k\pi \), \( k\in\mathbb{Z} \)}{}{}}
\End

\Data\ID{1103055206}
\Updated{2021-01-13 01:01:20}
\Level{B}
\SubArea{Angles, arcs and sectors}
\IMG{1103055206.pdf}
\LANGen{
\Question{
\( ABCD \) is a square, as it is shown in the picture. Find the measures of all the coterminal angles to the angle \( BDA \). (Two angles are coterminal if they are drawn in the standard position and both have their terminal sides in the same location.)}
{\( \frac74\pi+2k\pi \), \( k\in\mathbb{Z} \)}{\( \frac4\pi+2k\pi \), \( k\in\mathbb{Z} \)}{\( \frac4\pi+k\pi \), \( k\in\mathbb{Z} \)}{\( -\frac74\pi+2k\pi \), \( k\in\mathbb{Z} \)}{}{}}
\LANGpl{
\Question{
Dany jest kwadrat \( ABCD \). Wszystkie kąty skierowane na  \( BDA \) można zapisać jako:}
{\( \frac74\pi+2k\pi \), \( k\in\mathbb{Z} \)}{\( \frac4\pi+2k\pi \), \( k\in\mathbb{Z} \)}{\( \frac4\pi+k\pi \), \( k\in\mathbb{Z} \)}{\( -\frac74\pi+2k\pi \), \( k\in\mathbb{Z} \)}{}{}}
\LANGcs{
\Question{
Je dán čtverec \( ABCD \). Všechny velikosti orientovaného úhlu  \( BDA \) je možné zapsat ve tvaru:}
{\( \frac74\pi+2k\pi \), \( k\in\mathbb{Z} \)}{\( \frac4\pi+2k\pi \), \( k\in\mathbb{Z} \)}{\( \frac4\pi+k\pi \), \( k\in\mathbb{Z} \)}{\( -\frac74\pi+2k\pi \), \( k\in\mathbb{Z} \)}{}{}}
\LANGsk{
\Question{
Daný je štvorec \( ABCD \). Všetky veľkosti orientovaného uhla  \( BDA \) možno zapísať v tvare:}
{\( \frac74\pi+2k\pi \), \( k\in\mathbb{Z} \)}{\( \frac4\pi+2k\pi \), \( k\in\mathbb{Z} \)}{\( \frac4\pi+k\pi \), \( k\in\mathbb{Z} \)}{\( -\frac74\pi+2k\pi \), \( k\in\mathbb{Z} \)}{}{}}
\LANGes{
\Question{
Sea un cuadrado \( ABCD \). Todas las medidas del ángulo orientado \( BDA \) se pueden escribir de forma:}
{\( \frac74\pi+2k\pi \), \( k\in\mathbb{Z} \)}{\( \frac4\pi+2k\pi \), \( k\in\mathbb{Z} \)}{\( \frac4\pi+k\pi \), \( k\in\mathbb{Z} \)}{\( -\frac74\pi+2k\pi \), \( k\in\mathbb{Z} \)}{}{}}
\End

\Data\ID{9000045710}
\Updated{2021-01-13 12:01:38}
\Level{B}
\SubArea{Angles, arcs and sectors}
\LANGen{
\Question{
Find the length \(l\) 
of a latitude at \(50^{\circ }\) 
N. (Use \(R\) 
for the radius of the Earth.)}
{\(l = 2\pi R\cos 50^{\circ }\)}{\(l = 2\pi R\sin 50^{\circ }\)}{\(l = 2\pi R\mathop{\mathrm{tg}}\nolimits 50^{\circ }\)}{\(l = 2\pi R\mathop{\mathrm{cotg}}\nolimits 50^{\circ }\)}{}{}}
\LANGpl{
\Question{
Oblicz długość  \(l\) 
szerokości geograficznej  \(50^{\circ }\) północ. (\(R\) 
to promień Ziemi.)}
{\(l = 2\pi R\cos 50^{\circ }\)}{\(l = 2\pi R\sin 50^{\circ }\)}{\(l = 2\pi R\mathop{\mathrm{tg}}\nolimits 50^{\circ }\)}{\(l = 2\pi R\mathop{\mathrm{cotg}}\nolimits 50^{\circ }\)}{}{}}
\LANGcs{
\Question{
Určete vztah, který platí pro délku
\(l\) rovnoběžky na
\(50^{\circ }\) severní šířce.
(Symbolem \(R_{Z}\) 
značíme poloměr Země.)}
{\(l = 2\pi R_{Z}\cos 50^{\circ }\)}{\(l = 2\pi R_{Z}\sin 50^{\circ }\)}{\(l = 2\pi R_{Z}\mathop{\mathrm{tg}}\nolimits 50^{\circ }\)}{\(l = 2\pi R_{Z}\mathop{\mathrm{cotg}}\nolimits 50^{\circ }\)}{}{}}
\LANGsk{
\Question{
Určte vzťah, ktorý platí pre dĺžku
\(l\) rovnobežky na
\(50^{\circ }\) severnej šírke.
(Symbolom \(R_{Z}\) 
značíme polomer Zeme.)}
{\(l = 2\pi R_{Z}\cos 50^{\circ }\)}{\(l = 2\pi R_{Z}\sin 50^{\circ }\)}{\(l = 2\pi R_{Z}\mathop{\mathrm{tg}}\nolimits 50^{\circ }\)}{\(l = 2\pi R_{Z}\mathop{\mathrm{cotg}}\nolimits 50^{\circ }\)}{}{}}
\LANGes{
\Question{
Averigua la relación que vale para la longitud 
\(l\) de la paralela en 
\(50^{\circ }\) de longitud norte 
(El símbolo \(R_{T}\) 
el para el radio de la Tierra.)}
{\(l = 2\pi R_{T}\cos 50^{\circ }\)}{\(l = 2\pi R_{T}\sin 50^{\circ }\)}{\(l = 2\pi R_{T}\mathop{\mathrm{tg}}\nolimits 50^{\circ }\)}{\(l = 2\pi R_{T}\mathop{\mathrm{cotg}}\nolimits 50^{\circ }\)}{}{}}
\End
